% !TEX root = sga4.5-en.tex
% ENGLISH EDITION -- TRANSLATED FROM FRENCH (AI-assisted, needs human review).

\chapter{\'Etale cohomology: the starting points}\label{I}
\blfootnote{by P.\ Deligne, notes by J.F.\ Boutot}




















\section{Grothendieck topologies}\label{I:1}

Originally, Grothendieck topologies appeared as underlying
his theory of descent (cf. SGA 1 VI, VIII); the use of the corresponding
cohomology theories came later. The same approach is followed
here: by formalizing the classical notions of localization, local
property and gluing (\ref{I:1-1}, \ref{I:1-2}, \ref{I:1-3}), one arrives at
the general concept of Grothendieck topology (\ref{I:1-6}); to justify
its introduction in algebraic geometry, we prove a faithfully
flat descent theorem (\ref{I:1-4}), generalizing the classical Hilbert
theorem (\ref{I:1-5}).

The reader will find a more complete, yet concise, exposition of the
formalism in Giraud \cite{gi64}. M.\ Artin's notes: ``Grothendieck
topologies'' \cite{ar62} (chapters I to III) remain useful as well. The
866 pages of expos\'es I to V of SGA 4 are valuable when
considering exotic topologies, such as the one giving rise to
crystalline cohomology; to use the \'etale topology, so close to
classical intuition, it is not indispensable to read them. 










\subsection{Sieves}\label{I:1-1}

Let $X$ be a topological space and $f:X\to \dR$ a real-valued
function on $X$. Continuity of $f$ is a property of a
local nature; in other words, if $f$ is continuous on every sufficiently small
open subset of $X$, then $f$ is continuous on all of $X$. To formalize the notion of
``property of a local nature,'' we introduce a few definitions.

We say that a set $\cU$ of open subsets of $X$ is a \emph{sieve} if for every
$U\in\cU$ and $V\subset U$, we have $V\in\cU$. We say that a sieve is
\emph{covering} if the union of all the open subsets belonging to this sieve is
equal to $X$.

Given a family $\{U_i\}$ of open subsets of $X$, the sieve generated by
$\{U_i\}$ is by definition the set of open subsets $U$ of $X$ such that $U$
is contained in one of the $U_i$.

We say that a property $P(U)$, defined for every open subset $U$ of $X$, is
\emph{local} if, for every covering sieve $\cU$ of every open subset $U$ of $X$,
$P(U)$ holds if and only if $P(V)$ holds for every $V\in \cU$. For
example, given $f:X\to \dR$, the property ``$f$ is continuous on $U$''
is local. 





\subsection{Sheaves}\label{I:1-2}

Let us make precise the notion of a function given locally on $X$.





\subsubsection{Sieve point of view}\label{I:1-2-1}

Let $\cU$ be a sieve of open subsets of $X$. By a function given
$\cU$-locally on $X$ we mean the data, for every $U\in \cU$, of a function $f_U$
on $U$ such that, if $V\subset U$, we have $f_V=f_U|V$. 





\subsubsection{\v Cech point of view}\label{I:1-2-2}

If the sieve $\cU$ is generated by a family of open subsets $U_i$ of $X$, to
give a function $\cU$-locally amounts to giving a function $f_i$
on each $U_i$, such that $f_i|{U_i\cap U_j} = f_j|{U_i\cap U_j}$.

In other words, if $Z=\coprod U_i$, to give a function $\cU$-locally
amounts to giving a function on $Z$ which is constant on the fibres of
the natural projection $Z\to X$. 





\subsubsection{}\label{I:1-2-3}

Continuous functions form a sheaf; this means that for every covering
sieve $\cU$ of an open subset $V$ of $X$ and every function given
$\cU$-locally $\{f_U\}$ such that each $f_U$ is continuous on $U$, there
exists a unique continuous function $f$ on $V$ such that $f|U=f_U$ for every
$U\in \cU$.










\subsection{Stacks}\label{I:1-3}

Let us now make precise the notion of a vector bundle given locally on $X$. 





\subsubsection{Sieve point of view}\label{I:1-3-1}

Let $\cU$ be a sieve of open subsets of $X$. By a vector bundle given
$\cU$-locally on $X$ we mean the data of 
\begin{enumerate}[\indent a)]
  \item a vector bundle $E_U$ on each $U\in \cU$, 
  \item if $V\subset U$, an isomorphism $\rho_{U,V} : E_V\iso E_U|V$, satisfying 
  \item if $W\subset V\subset U$, the diagram 
    \[\xymatrix{
      E_W \ar[r]^-{\rho_{U,W}} \ar[dr]_-{\rho_{V,W}}
        & E_U|W \\
      & E_V|W \ar[u]_-{\rho_{U,V}|W}
    }\]
    commute, i.e. 
    $\rho_{U,V} = (\mbox{$\rho_{U,V}$ restricted to $W$}) \circ \rho_{V,W}$. 
\end{enumerate}





\subsubsection{\v Cech point of view}\label{I:1-3-2}

If the sieve $\cU$ is generated by a family of open subsets $U_i$ of $X$, to
give a vector bundle $\cU$-locally amounts to giving:
\begin{enumerate}[\indent a)]
  \item a vector bundle $E_i$ on each $U_i$, 
  \item if $U_{ij} = U_i\cap U_j = U_i\times_X U_j$, an isomorphism 
    $\rho_{ji} : E_i|U_{ij} \iso E_j|U_{ij}$, such that 
  \item if $U_{ijk} = U_i\times_X U_j\times_X U_k$, the diagram 
    \[\xymatrix{
      E_i |U_{ijk} \ar[r]^-{\rho_{ki}|U_{ijk}} \ar[dr]_-{\rho_{ji}|U_{ijk}} 
        & E_k|U_{ijk} \\
      & E_j|U_{ijk} \ar[u]_-{\rho_{kj}|U_{ijk}}
    }\]
    commute, i.e. $\rho_{ki} = \rho_{kj}\circ \rho_{ji}$ on
    $U_{ijk}$. 
\end{enumerate}

In other words, if $Z=\coprod U_i$ and $\pi:Z\to X$ is the natural
projection, to give a vector bundle $\cU$-locally amounts to giving: 
\begin{enumerate}[\indent a)]
  \item a vector bundle $E$ on $Z$, 
  \item if $x$ and $y$ are two points of $Z$ such that $\pi(x) = \pi(y)$, an
    isomorphism $\rho_{y x}:E_x\iso E_y$ between the fibres of $E$ at $x$ and at
    $y$, depending continuously on $(x,y)$ and such that, 
  \item if $x$, $y$ and $z$ are three points of $Z$ such that 
  $\pi(x)=\pi(y)=\pi(z)$, we have $\rho_{zx} = \rho_{zy}\circ \rho_{yx}$. 
\end{enumerate}





\subsubsection{}\label{I:1-3-3}

A vector bundle $E$ on $X$ defines a vector bundle given
$\cU$-locally $E_\cU$; the system of restrictions $E_U$ of $E$ to the objects
of $\cU$. The fact that the notion of vector bundle is of a local nature can
be expressed as follows: for every covering sieve $\cU$ of $X$, the functor
$E\mapsto E_\cU$, from vector bundles on $X$ to vector bundles
given $\cU$-locally, is an equivalence of categories. 





\subsubsection{}\label{I:1-3-4}

If in \ref{I:1} one replaces ``open subset of $X$'' by ``subset of
$X$,'' one obtains the notion of a sieve of subspaces of $X$. In this setting
too there are gluing theorems. For example: let $X$ be a
space and $\cC$ a sieve of subspaces of $X$ generated by a locally finite
closed covering of $X$, then the functor $E\mapsto E_\cC$, from vector
bundles on $X$ to vector bundles given $\cC$-locally, is an
equivalence of categories.

In algebraic geometry, it is useful to consider also ``sieves
of spaces over $X$''; this is what we shall see in the next paragraph. 










\subsection{Faithfully flat descent}\label{I:1-4}





\subsubsection{}\label{I:1-4-1}

In the setting of schemes, the Zariski topology is not fine enough for
the study of non-linear problems and one is led to replace, in the
preceding definitions, open immersions by more
general morphisms. From this point of view, descent techniques appear as
localization techniques. Thus the following descent statement can be expressed
by saying that the properties under consideration are of a local nature for the
faithfully flat topology (a morphism of schemes is said to be faithfully
flat if it is flat and surjective). 


\begin{proposition}\label{I:1-4-2}
Let $A$ be a ring and $B$ a faithfully flat $A$-algebra. Then:
\begin{enumerate}[(i)]
  \item A sequence $\Sigma=(M'\to M\to M'')$ of $A$-modules is exact as soon
    as the sequence $\Sigma_{(B)}$ obtained from it by scalar extension
    to $B$ is exact.
  \item An $A$-module $M$ is of finite type (resp. of finite presentation,
    flat, locally free of finite rank, invertible (i.e.\ locally free
    of rank one)) as soon as the $B$-module $M_{(B)}$ is.
\end{enumerate}
\end{proposition}
\begin{proof}
(i) Since the functor $M\mapsto M_{(B)}$ is exact (flatness of $B$), it suffices
to show that, if an $A$-module $N$ is non-zero, then $N_{(B)}$ is non-zero. If
$N$ is non-zero, $N$ contains a non-zero cyclic submodule $A/\fa$; then
$N_{(B)}$ contains a cyclic submodule $(A/\fa)_{(B)} = B/\fa B$, non-zero
by surjectivity of the structural morphism $\varphi:\spec(B)\to\spec(A)$ (if
$V(\fa)$ is non-empty, $\varphi^{-1}(V(\fa))=V(\fa B)$ is non-empty).

(ii) For any family $(x_i)$ of elements of $M_{(B)}$, there exists a
finitely generated submodule $M'$ of $M$ such that $M'_{(B)}$ contains the $x_i$. If
$M_{(B)}$ is of finite type and the $x_i$ generate $M_{(B)}$, we have
$M'_{(B)}=M_{(B)}$, hence $M'=M$ and $M$ is of finite type.
\end{proof}

If $M_{(B)}$ is of finite presentation, one can, by what precedes,
find a surjection $A^n\to M$. If $N$ is the kernel of this surjection, the
$B$-module $N_{(B)}$ is of finite type, hence so is $N$, and $M$ is of
finite presentation. The assertion for ``flat'' follows at once from (i);
``locally free of finite rank'' means ``flat and of finite presentation''
and the rank is tested by scalar extension to fields.





\subsubsection{}\label{I:1-4-3}

Let $X$ be a scheme and $\cS$ a class of $X$-schemes stable under fibre
product over $X$. A class $\cU\subset \cS$ is a \emph{sieve} over $X$
(relative to $\cS$) if, for every morphism $\varphi:V\to U$ of
$X$-schemes, with $U,V\in \cS$ and $U\in\cU$, we have $V\in \cU$. The sieve
\emph{generated} by a family $\{U_i\}$ of $X$-schemes in $\cS$ is the
class of $V\in\cS$ such that there exists a morphism of $X$-schemes from $V$ to
one of the $U_i$. 






\subsubsection{}\label{I:1-4-4}

Let $\cU$ be a sieve over $X$. By a quasi-coherent module given
$\cU$-locally on $X$ we mean the data of
\begin{enumerate}[\indent a)]
  \item a quasi-coherent module $E_U$ on each $U\in\cU$, 
  \item for every $U\in\cU$ and every morphism $\varphi:V\to U$ of
    $X$-schemes in $\cS$, an isomorphism
    $\rho_\varphi:E_V\iso \varphi^* E_U$, these being such that
  \item if $\psi:W\to V$ is a morphism of $X$-schemes in $\cS$, the
    diagram 
    \[\xymatrix{
      E_W \ar[r]^-{\rho_{\varphi\circ\psi}} \ar[dr]_-{\rho_\psi} 
        & \psi^*\varphi^* E_U \\
      & \psi^* E_V \ar[u]_-{\psi^*\rho_\varphi}
    }\]
    commute, i.e.
    $\rho_{\varphi\circ\psi}=(\psi^*\rho_\varphi)\circ \rho_\psi$. 
\end{enumerate}

If $E$ is a quasi-coherent module on $X$, we denote by $E_\cU$ the module
given $\cU$-locally with value $\varphi_U^* E$ on $\varphi:U\to X$ and such that,
for every morphism $\psi:V\to U$ the restriction isomorphism $\rho_\psi$
is the canonical isomorphism
$E_V=(\varphi_U\circ \psi)^* E\iso \psi^* \varphi_U^* E = \psi^* E_U$. 


\begin{theorem}\label{I:1-4-5}
Let $\{U_i\}\in\cS$ be a finite family of $X$-schemes flat over $X$ such that
$X$ is the union of the images of the $U_i$, and let $\cU$ be the sieve generated by
$\{U_i\}$. Then the functor $E\mapsto E_\cU$ is an equivalence from the
category of quasi-coherent modules on $X$ to the category of quasi-coherent
modules given $\cU$-locally.
\end{theorem}
\begin{proof}
We shall treat only the case where $x$ is affine and $\cU$ is generated by
a $X$-scheme $U$ affine faithfully flat over $X$. The reduction to this case is
formal. Put $X=\spec(A)$ and $U=\spec(B)$.

If the morphism $U\to X$ admits a section, $X$ belongs to the sieve $\cU$ and
the assertion is clear. We shall reduce to this case.

A quasi-coherent module given $\cU$-locally defines modules $M'$, $M''$
and $M'''$ over $U$, $U\times_X U$ and $U\times_X U\times_X U$, and
isomorphisms $\rho:p^* M^\bullet\simeq M^\bullet$ for every
projection morphism $p$ between these spaces; this is a \emph{cartesian diagram} 
\[\xymatrix{
  M^* : M' \ar@<2pt>[r] \ar@<-2pt>[r] 
    & M'' \ar[r] \ar@<4pt>[r] \ar@<-4pt>[r] 
    & M'''
}\]
above
\[\xymatrix{
  U_* : U 
    & U\times_X U \ar@<2pt>[l] \ar@<-2pt>[l] 
    & U\times_X U\times_X U \ar[l] \ar@<4pt>[l] \ar@<-4pt>[l] \mbox{.}
}\]
Conversely $M^*$ determines the module given $\cU$-locally: for
$V\in\cU$, there exists $\varphi:V\to U$ and we put $M_U=\varphi^* M'$;
for $\varphi_1,\varphi_2:V\to U$, we have a natural identification
$\varphi_1^* M'\simeq (\varphi_1\times \varphi_2)^* M'' \simeq \varphi_2^* M'$,
and one sees using $M'''$ that these identifications are compatible, so
that the definition is legitimate. In short, giving a
module $\cU$-locally amounts to the same as giving a cartesian diagram $M^*$ over $U_*$.

In algebraic terms: giving $M^*$ amounts to giving a cartesian
diagram of modules 
\[\xymatrix{
  M' \ar@<3pt>[r]|-{\partial_0} \ar@<-3pt>[r]|-{\partial_1} 
    & M'' \ar@<6pt>[r]|-{\partial_0} \ar[r]|-{\partial_1} \ar@<-6pt>[r]|-{\partial_2}
    & M'''
}\]
above the diagram of rings 
\[\xymatrix{
  B \ar@<3pt>[r]|-{\partial_0} \ar@<-3pt>[r]|-{\partial_1} 
    & B\otimes_A B \ar@<6pt>[r]|-{\partial_0} \ar[r]|-{\partial_1} \ar@<-6pt>[r]|-{\partial_2} 
    & B\otimes_A B\otimes_A B
}\]
(more precisely: we have $\partial_i(b m)=\partial_i(b)\cdot\partial_i(m)$, the
usual identities such as $\partial_0\partial_1=\partial_0\partial_0$
hold, and ``cartesian'' means that the morphisms
$\partial_i:M'\otimes_{B,\partial_i}(B\otimes_A B)\to M''$ and
$M''\otimes_{B\otimes_A B,\partial_i}(B\otimes_A B\otimes_A B)\to M'''$ are
isomorphisms).

The functor $E\mapsto E_\cU$ becomes the functor which, to an $A$-module $M$,
associates 
\[\xymatrix{
  M^* = (M\otimes_A B \ar@<2pt>[r] \ar@<-2pt>[r] 
    & M\otimes_A B\otimes_A B \ar@<-4pt>[r] \ar[r] \ar@<4pt>[r] 
    & M\otimes_A B\otimes_A B\otimes_A B)\text{.}
}\]
It admits as right adjoint the functor
\[\xymatrix{
  (M' \ar@<2pt>[r] \ar@<-2pt>[r] 
    & M'' \ar@<-4pt>[r] \ar[r] \ar@<4pt>[r]
    & M''') \ar@{|->}[r] 
    & \ker(M' \ar@<2pt>[r] \ar@<-2pt>[r] 
    & M''')\text{.}
}\]
We must prove that the adjunction maps 
\[
  M \to\ker(M\otimes_A B\rightrightarrows M\otimes_A B\otimes_A B)
\]
et
\[
  \ker(M'\rightrightarrows M'')\otimes_A B \to M'
\]
are isomorphisms. By (\ref{I:1-4-2}.i), it suffices to prove this
after a faithfully flat base change $A\to A'$ ($B$ becoming
$B'=B\otimes_A A'$). Taking $A'=B$, this reduces us to the case where $U\to X$
admits a section.
\end{proof}










\subsection{A special case: Hilbert's theorem 90}\label{I:1-5}





\subsubsection{}\label{I:1-5-1}

Let $k$ be a field, $k'$ a Galois extension of $k$ and $G=\gal(k'/k)$.
Then the homo\~morphism 
\begin{align*}
  k'\otimes_k k' &\to \oplus_{\sigma\in G} k' \\
  x\otimes y     &\mapsto \{x\cdot \sigma(y)\}_{\sigma\in G}
\end{align*}
is bijective. 

It follows that giving a module locally for 
the sieve generated by $\spec(k')$ over $\spec(k)$ amounts to the same as giving a 
$k'$-vector space equipped with a semi-linear action of $G$, i.e.: 
\begin{enumerate}[\indent a)]
  \item a $k'$-vector space $V'$,
  \item for every $\sigma\in G$, an endomorphism $\varphi_\sigma$ of the 
    group structure of $V'$ such that 
    $\varphi_\sigma(\lambda v) = \sigma(\lambda) \varphi_\sigma(v)$, for
    all $\lambda\in k'$ and $v\in V'$, satisfying the condition
  \item for all $\sigma,\tau\in G$, we have 
    $\varphi_{\tau\sigma} = \varphi_\tau\circ\varphi_\sigma$. 
\end{enumerate}

Let $V={V'}^G$ be the group of invariants under this action of $G$; it is a 
$k$-vector space and, by Theorem \ref{I:1-4-5}, we have:

\begin{proposition}\label{I:1-5-2}
The inclusion of $V$ into $V'$ defines an isomorphism 
$V\otimes_k k' \iso V'$.
\end{proposition}

In particular, if $V'$ is of dimension $1$ and $v'\in V$ is non-zero, 
$\varphi_\sigma$ is determined by the constant 
$c(\sigma)\in {k'}^\times$ such that $\varphi_\sigma(v')=c(\sigma)v'$ and 
condition c) reads 
\[
  c(\tau\sigma) = c(\tau)\cdot \tau(c(\sigma))\text{.}
\]
By the proposition there exists a non-zero invariant vector 
$v=\mu v'$, $\mu\in {k'}^\times$. Hence for every $\sigma\in G$, 
\[
  c(\sigma) = \mu\cdot \sigma(\mu^{-1})\text{.}
\]
In other words every $1$-cocycle of $G$ with values in ${k'}^\times$ is a 
coboundary: 

\begin{corollary}\label{I:1-5-3}
We have $\h^1(G,{k'}^\times)=0$.
\end{corollary}










\subsection{Grothendieck topologies}\label{I:1-6}

We now transcribe the definitions of the preceding paragraphs 
into an abstract setting encompassing both the case of topological spaces 
and that of schemes. 





\subsubsection{}\label{I:1-6-1}

Let $\cS$ be a category and $U$ an object of $\cS$. By a \emph{sieve} 
over $U$ we mean a subset $\cU$ of $\ob(\cS/U)$ such that if $\varphi:V\to U$ 
belongs to $\cU$ and $\psi:W\to V$ is a morphism in $\cS$, then 
$\varphi\circ\psi:W\to U$ belongs to $\cU$. 

If $\{\varphi_i:U_i\to U\}$ is a family of morphisms, the sieve 
generated by the $U_i$ is by definition the set of morphisms 
$\varphi:V\to U$ which factor through one of the $\varphi_i$. 

If $\cU$ is a sieve over $U$ and $\varphi:V\to U$ is a morphism, the 
restriction $\cU_V$ of $\cU$ to $V$ is by definition the sieve over $V$ 
consisting of the morphisms $\psi:w\to V$ such that 
$\varphi\circ\psi:W\to U$ belongs to $\cU$.





\subsubsection{}\label{I:1-6-2}

The data of a \emph{Grothendieck topology} on $\cS$ consists of the data, for every object $U$ of $\cS$, 
of a set $C(U)$ of sieves over $U$, 
called covering sieves, such that the following axioms hold: 
satisfaits:
\begin{enumerate}[\indent a)]
  \item The sieve generated by the identity of $U$ is covering.
  \item If $\cU$ is a covering sieve over $U$ and $V\to U$ is a 
    morphism, the sieve $\cU_V$ is covering.
  \item A locally covering sieve is covering. In other words, if $\cU$ is 
    a covering sieve over $U$ and $\cU'$ is a sieve over $U$ such that, for 
    every $V\to U$ belonging to $\cU$, the sieve $\cU_V'$ is covering, then 
    $\cU'$ is covering.
\end{enumerate}

We call \emph{site} the data of a category equipped with a Grothendieck topology. 
Grothendieck.





\subsubsection{}\label{I:1-6-3}

Given a site $\cS$, by a presheaf on $\cS$ we mean a contravariant functor 
$\sF$ from $\cS$ to the category of sets. For every 
object $U$ of $\cS$, by a section of $\sF$ over $U$ we mean the 
elements of $\sF(U)$. For every morphism $V\to U$ and every 
$s\in \sF(U)$, we denote by $s|V$ ($s$ restricted to $V$) the image of $s$ in 
$\sF(V)$.

If $\cU$ is a sieve over $U$, by a section given $\cU$-locally we mean the data, 
for every $V\to U$ belonging to $\cU$, of a section 
$s_V\in\sF(V)$ such that, for every morphism $W\to V$, we have $s_V|W=s_W$. We 
say that $\sF$ is a \emph{sheaf} if, for every object $U$ of $\cS$, for 
every covering sieve $\cU$ over $U$ and every section given 
$\cU$-locally $\{s_V\}$, there exists a unique section $s\in\sF(U)$ such that 
$s|V=s_V$, for every $V\to U$ belonging to $\cU$. 

Abelian sheaves are defined analogously by 
replacing the category of sets by that of abelian groups. One 
shows that the category of abelian sheaves on $\cS$ is an 
abelian category with enough injectives. A sequence 
$\sF\xrightarrow f \sG\xrightarrow g \sH$ of sheaves is exact if, for 
every object $U$ of $\cS$, and every $s\in \sG(U)$ such that $g(s)=0$, 
there locally exists $t$ such that $f(t)=s$; i.e.\ if there exists a covering sieve 
$\cU$ over $U$ and for every $V\in\cU$, a section $t_V$ of $\sF$ over $V$ such that 
that $f(t_V)=s|V$. 





\subsubsection{Examples}\label{I:1-6-4}

We have seen two of them above. 

\begin{enumerate}[\indent a)]
  \item Let $X$ be a topological space and $\cS$ the category whose objects 
    are the open subsets of $X$ and whose morphisms are the natural inclusions. The 
    Grothendieck topology on $\cS$ corresponding to the usual topology 
    of $X$ is the one for which a sieve $\cU$ over an open subset $U$ of $X$ is 
    covering if the union of the open subsets belonging to this sieve equals 
    $U$. It is clear that the category of sheaves on $\cS$ is equivalent 
    to the category of sheaves on $X$ in the usual sense. 
  \item Let $X$ be a scheme and $\cS$ the category of schemes over $X$. We 
    call fpqc topology (faithfully flat quasi-compact) on $\cS$ the 
    Grothendieck topology for which a sieve over an $X$-scheme $U$ is 
    covering if it is generated by a finite family of flat morphisms whose  
    images cover $U$. 
\end{enumerate}





\subsubsection{Cohomology}\label{I:1-6-5}

We shall always assume that the category $\cS$ has a final object $X$. Then we 
call global sections of an abelian sheaf $\sF$, and denote by 
$\Gamma\sF$ or $\h^0(X,\sF)$, the group $\sF(X)$. The functor 
$\sF\mapsto\Gamma\sF$ is a left exact functor from the category of 
abelian sheaves on $\cS$ to the category of abelian groups; we denote by 
$\h^i(X,-)$ its derived functors (or satellites). These cohomology groups 
represent the obstructions to passing from the local to the global. By definition, if 
$0\to\sF\to\sG\to\sH\to 0$ is an exact sequence of abelian sheaves, we have a long 
exact cohomology sequence:
\[\xymatrix{
  0 \ar[r] 
    & \h^0(X,\sF) \ar[r] 
    & \h^0(X,\sG) \ar[r] 
    & \h^0(X,\sH) \ar[r] 
    & \h^1(X,\sF) \ar[r] 
    & \cdots \\
  \ldots \ar[r] 
    & \h^n(X,\sF) \ar[r] 
    & \h^n(X,\sG) \ar[r]
    & \h^n(X,\sH) \ar[r] 
    & \h^{n+1}(X,\sF) \ar[r] 
    & \cdots
}\]





\subsubsection{}\label{I:1-6-6}

Given an abelian sheaf $\sF$ on $\cS$, we call 
\emph{$\sF$-torsor} a sheaf $\sG$ equipped with an action $\sF\times\sG\to\sG$ 
of $\sF$ such that locally (after restriction to all objects of a
covering sieve of the final object $X$) $\sG$ with the action of $\sF$ is 
isomorphic to $\sF$ equipped with the canonical action $\sF\times \sF\to \sF$ by 
translations. 

One can show that $\h^1(X,\sF)$ is interpreted as the set of isomorphism classes 
of $\sF$-torsors. 




















\section{\'Etale topology}\label{I:2}

We specialize the definitions of the preceding chapter to the case of the 
\'etale topology of a scheme $X$ (\ref{I:2-1}, \ref{I:2-2}, \ref{I:2-3}). The 
corresponding cohomology coincides in the case where $X$ is the spectrum of a 
field $K$ with the Galois cohomology of $K$ (\ref{I:2-4}). 










\subsection{\'Etale topology}\label{I:2-1}

We begin with some reminders on the notion of \'etale morphism. 

\begin{definition}\label{I:2-1-1}
Let $A$ be a (commutative) ring. An $A$-algebra $B$ is said to be \'etale if 
$B$ is an $A$-algebra of finite presentation and the following 
equivalent conditions hold:
\begin{enumerate}[\indent a)]
  \item For every $A$-algebra $C$ and every square-zero ideal $J$ of 
    $C$, the canonical map 
    \[
      \hom_{A\textnormal{-alg}}(B,C) \to \hom_{A\textnormal{-alg}}(B,C/J)
    \]
    is a bijection.
  \item $B$ is a flat $A$-module and $\Omega_{B/A}=0$ (we denote by $\Omega_{B/A}$ 
    the module of relative differentials).
  \item Let $B=A[X_1,\dotsc,X_n]/I$ be a presentation of $B$. Then for every 
    prime ideal $\fp$ of $A[X_1,\dotsc,X_n]$ containing $I$, there exist 
    polynomials $P_1,\dotsc,P_n\in I$ such that $I_\fp$ is generated by the 
    images of $P_1,\dotsc,P_n$ and $\det(\partial P_i/\partial X_j)\notin \fp$.  
\end{enumerate}
\end{definition}
(cf. \cite[I]{sga1} or \cite[V]{ra70})
%[cf. SGA I, exposé I ou M. {\sc Raynoud}, \emph{Anneaux Locaux Henséliens}, chapitre V]. 

A morphism of schemes $f:X\to S$ is said to be \emph{\'etale} if for every 
$x\in X$ there exist an affine open neighbourhood $U=\spec(A)$ of $f(x)$ and an 
affine open neighbourhood $V=\spec(B)$ of $x$ in $X\times_S U$ such that $B$ is 
an \'etale $A$-algebra. 





\subsubsection{Examples}\label{I:2-1-2}
\begin{enumerate}[\indent a)]
  \item If $A$ is a field, an $A$-algebra $B$ is \'etale if and only if 
    it is a finite product of separable extensions of $A$. 
  \item If $X$ and $S$ are schemes of finite type over $\dC$, a morphism 
    $f:X\to S$ is \'etale if and only if its analytification 
    $f^{\text{an}}:X^{\text{an}}\to Y^{\text{an}}$ is a local isomorphism.
\end{enumerate}





\subsubsection{Formal properties}\label{I:2-1-3}
\begin{enumerate}[\indent a)]
  \item (base change) If $f:X\to S$ is an \'etale morphism, so is 
    $f_{S'}:X\times_S S'\to S'$ for every morphism $S'\to S$. 
  \item If $f:X\to S$ and $g:Y\to S$ are two \'etale morphisms, every 
    $S$-morphism from $X$ to $Y$ is \'etale.
  \item (descent) Let $f:X\to S$ be a morphism. If there exists a 
    faithfully flat morphism $S'\to S$ such that $f_{S'}:X\times_S S'\to S'$ is 
    \'etale, then $f$ is \'etale.
\end{enumerate}





\subsubsection{}\label{I:2-1-4}

Let $X$ be a scheme. Let $\cS$ be the category of \'etale $X$-schemes; 
by (\ref{I:2-1-3}.c) every morphism of $\cS$ is an \'etale morphism. 
We call \emph{\'etale topology} on $\cS$ the topology for which a 
sieve over $U$ is covering if it is generated by a finite family of 
morphisms $\varphi_i:U_i\to U$ such that the union of the images of the 
$\varphi_i$ covers $U$. We call \emph{\'etale site} of $X$, and denote by 
$\et X$, the site defined by $\cS$ with the \'etale topology. 










\subsection{Examples of sheaves}\label{I:2-2}





\subsubsection{Constant sheaf}\label{I:2-2-1}

Let $C$ be an abelian group and assume for simplicity that $X$ is noetherian. We 
denote by $\const C_X$ (or even $C$ if there is no ambiguity) the sheaf 
defined by $U\mapsto C^{\pi_0(U)}$, where $\pi_0(U)$ is the (finite) set of 
connected components of $U$. The most important case will be $C=\dZ/n$. Thus by definition 
 
\[
  \h^0(X,\dZ/n)=\left(\dZ/n\right)^{\pi_0(X)}\text{.}
\]
Moreover $\h^1(X,\dZ/n)$ is the set of isomorphism classes of 
$\dZ/n$-torsors (\ref{I:1-6-6}), i.e.\ of Galois \'etale coverings 
of $X$ with group $\dZ/n$. In particular, if $X$ is connected and 
$\pi_1(X)$ is its fundamental group for a chosen base point, we have 
\[
  \h^1(X,\dZ/n) = \hom(\pi_1(X),\dZ/n)\text{.}
\]





\subsubsection{Multiplicative group}\label{I:2-2-2}

We denote by $\dG_{m,X}$ (or $\dG_m$ if there is no ambiguity) the sheaf 
defined by $U\mapsto \Gamma(U,\sO_U^\times)$; it is indeed a sheaf 
thanks to the faithfully flat descent theorem (\ref{I:1-4-5}). We have 
 
\[
  \h^0(X,\dG_m) = \h^0(X,\sO_X)^\times\text{;}
\]
in particular if $X$ is reduced, connected and proper over an 
algebraically closed field $k$, we have:
\[
  \h^0(X,\dG_m) = k^\times\text{.}
\]

\begin{proposition}\label{I:2-2-3}
We have an isomorphism:
\[
  \h^1(X,\dG_m) = \pic(X)\text{,}
\]
where $\pic(X)$ is the group of isomorphism classes of invertible sheaves on $X$.
\end{proposition}
\begin{proof}
Let $*$ be the functor which, to an invertible sheaf $\sL$ on $X$, associates the 
following presheaf $\sL^*$ on $\et X$: for $\varphi:U\to X$ \'etale, 
\[
  \sL^*(U)=\operatorname{Isom}_U(\sO_U,\varphi^*\sL)\text{.}
\]
% the original references here are just to (4.2) and (4.5), but the context 
% makes it seem that they refer to the results in chapter 1. Likewise a little 
% later
By (\ref{I:1-4-2}.i) and (\ref{I:1-4-5}) (full faithfulness), this 
presheaf is a sheaf; it is even a $\dG_m$-torsor. One checks 
at once that 
\begin{enumerate}[\indent a)]
  \item the functor $*$ is compatible with (\'etale) localization;
  \item it induces an equivalence from the category of 
    trivial invertible sheaves (i.e.\ isomorphic to $\sO_X$) to the category of 
    trivial $\dG_m$-torsors: $\sL$ is trivial if and only if $\sL^*$ 
    l'est.
\end{enumerate}

Moreover, by (\ref{I:1-4-2}.ii) and (\ref{I:1-4-5}), 
\begin{enumerate}[\indent a)]
\setcounter{enumi}{2}
  \item the notion of invertible sheaf is local for the \'etale topology. 
\end{enumerate}

It follows formally from a), b), c) that $*$ is an equivalence between the 
category of invertible sheaves on $X$ and that of $\dG_m$-torsors 
on $\et X$; it induces the desired isomorphism. The inverse 
equivalence is constructed as follows: if $\sT$ is a $\dG_m$-torsor, there exists a 
finite \'etale covering $\{U_i\}$ of $X$ such that the torsors $\sT/U_i$ are 
trivial; $\sT$ is then trivial over each $V$ \'etale over $X$ belonging to the 
sieve $\cU\subset \et X$ generated by $\{U_i\}$. Over each 
$V\in\cU$, $\sT|V$ corresponds to an invertible sheaf $\sL_V$ (by b)) and the 
$\sL_V$ form an invertible sheaf given $\cU$-locally $\sL_\cU$ 
(by a)). By c), the latter comes from an invertible sheaf $\sL(\sT)$ on 
$X$, and $\sT\mapsto \sL(\sT)$ is the desired inverse of $*$. 
\end{proof}





\subsubsection{Roots of unity}\label{I:2-2-4}

For every integer $n>0$, we call the sheaf of $n$-th roots of 
unity, and denote by $\dmu_n$, the kernel of raising to the 
$n$-th power in $\dG_m$. If $X$ is a scheme over a separably 
closed field $k$ and $n$ is invertible in $k$, the choice of a primitive 
$n$-th root of unity $\zeta\in k$ defines an isomorphism 
$i\mapsto \zeta^i$ from $\dZ/n$ onto $\dmu_n$. 

The relation between cohomology with coefficients in $\dmu_m$ and cohomology with 
coefficients in $\dG_m$ is given by the cohomology exact sequence deduced from the Kummer exact sequence below. 






\begin{theorem}[Kummer theory]\label{I:2-2-5}
If $n$ is invertible on $X$, raising to the $n$-th power in 
$\dG_m$ is an epimorphism of sheaves. Hence we have an exact sequence 
\[
  0 \to \dmu_n \to \dG_m \to \dG_m \to 0\text{.}
\]
\end{theorem}
\begin{proof}
Let $U\to X$ be an \'etale morphism and $a\in \dG_m(U) =\Gamma(U,\sO_U^\times)$. 
Since $n$ is invertible on $U$, the equation $T^n-a = 0$ is separable; 
in other words $U'=\spec\left(\sO_U[T]/(T^n-a)\right)$ is \'etale over 
$U$. Moreover $U'\to U$ is surjective and $a$ admits an $n$-th root over 
$U'$, whence the result. 
\end{proof}










\subsection{Fibres, direct images}\label{I:2-3}





\subsubsection{}\label{I:2-3-1}

We call \emph{geometric point} of $X$ a morphism $\bar x\to X$, where 
$\bar x$ is the spectrum of a separably closed field $k(\bar x)$. We denote it 
abusively by $\bar x$, the morphism $\bar x\to X$ being understood. If $x$ is 
the image of $\bar x$ in $X$, we say that $\bar x$ is centred at $x$. If the 
field $k(\bar x)$ is an algebraic extension of the residue field $k(x)$, we say 
that $\bar x$ is an \emph{algebraic} geometric point of $X$. 

We call \emph{\'etale neighbourhood} of $\bar x$ a commutative diagram 
\[\xymatrix{
  & U \ar[d] \\
  \bar x \ar[ur] \ar[r] 
  & X\text{,}
}\]
where $U\to X$ is an \'etale morphism. 

The \emph{strict localisation} of $X$ at $\bar x$ is the ring 
$\sO_{X,\bar x} = \varinjlim \Gamma(U,\sO_U)$, the inductive limit being over the 
\'etale neighbourhoods of $\bar x$. It is a strictly henselian local ring whose 
residue field is the separable closure of the residue field $k(x)$ of $X$ at 
$x$ in $k(\bar x)$. It plays the role of local ring for the \'etale topology. 





\subsubsection{}\label{I:2-3-2}

Given a sheaf $\sF$ on $\et X$, we call \emph{fibre} of $\sF$ at
$\bar x$ the set (resp. group,\dots) $\sF_{\bar x}=\varinjlim \sF(U)$,
the inductive limit always being taken over the \'etale neighbourhoods of $X$.

For a homomorphism of sheaves $\sF\to \sG$ to be a
mono-/epi-/isomorphism it is necessary and sufficient that the same hold for the morphisms
$\sF_{\bar x}\to \sG_{\bar x}$ induced on the fibres at every geometric point
of $X$. If $X$ is of finite type over an algebraically closed field, it suffices
that this hold at the rational points of $X$. 





\subsubsection{}\label{I:2-3-3}

If $f:X\to Y$ is a morphism of schemes and $\sF$ a sheaf on $\et X$,
the \emph{direct image} $f_* \sF$ of $\sF$ by $f$ is the sheaf on $\et Y$
defined by $f_* \sF(V) = \sF(X\times_Y V)$ for every $V$ \'etale over $Y$. The functor
$f_*:\sh(\et X)\to \sh(\et Y)$ is left exact. Its
right derived functors $\eR^q f_*$ are called higher direct images. If
$\bar y$ is a geometric point of $Y$, we have
\[
  (\eR^q f_* \sF)_{\bar y} = \varinjlim \h^q(V\times_Y X,\sF)\text{,}
\]
the inductive limit being taken over the \'etale neighbourhoods $V$ of $\bar y$. 

Let $\sO_{Y,\bar y}$ be the strict localisation of $Y$ at $\bar y$,
$\widetilde Y=\spec(\sO_{Y,\bar y})$ and $\widetilde X=X\times_Y \widetilde Y$.
One can extend $\sF$ to $\et{\widetilde X}$ (this is a special case of the
general notion of inverse image) as follows: let
$\widetilde U$ be a scheme \'etale over $\widetilde X$, then there exist an
\'etale neighbourhood $V$ of $\bar y$ and a scheme $U$ \'etale over $X\times_Y V$ such
that $\widetilde U=U\times_V \widetilde Y$; we put
\[
  \sF(\widetilde U)=\varinjlim \sF\left(U\times_V V'\right)\text{,}
\]
the inductive limit being taken over the \'etale neighbourhoods $V'$ of $\bar y$ which
dominate $V$. With this definition, we have
\[
  \left(\eR^q f_* \sF\right)_{\bar y} = \h^q(\widetilde X,\sF)\text{.}
\]

The functor $f_*$ has a left adjoint $f^*$, the ``inverse image'' functor.
If $\bar x$ is a geometric point of $X$ and $f(\bar x)$ its
image in $Y$, we have $(f^* \sF)_{\bar x} = \sF_{f(\bar x)}$. This formula shows
that $f^*$ is an exact functor. The functor $f_*$ therefore sends an injective
sheaf to an injective sheaf, and the spectral sequence of the composite functor
$\Gamma\circ f_*$ (resp. $g_* f_*$) gives the 





\begin{theorem}[Leray spectral sequence]\label{I:2-3-4}
Let $\sF$ be an abelian sheaf on $\et X$ and $f:X\to Y$ a morphism of
schemes (resp. morphisms of schemes $X\xrightarrow f Y \xrightarrow g Z$).
We have a spectral sequence 
\begin{align*}
  E_2^{pq} &= \h^p(Y,\eR^q f_* \sF) \Rightarrow \h^{p+q}(X,\sF) \\
  \text{(resp.}\qquad E_2^{pq} &= \eR^p g_* \eR^q f_* \sF \Rightarrow \eR^{p+q}(gf)_* \sF\text{).}
\end{align*}
\end{theorem}





\begin{corollary}\label{I:2-3-5}
If $\eR^q f_* \sF=0$ for all $q>0$, we have $\h^p(Y,f_* \sF)=\h^p(X,\sF)$ (resp.
$\eR^p g_*(\sF_* \sF) = \eR^p(g f)_* \sF$) for all $p\geqslant 0$.
\end{corollary}

This applies in particular in the following case:





\begin{proposition}\label{I:2-3-6}
Let $f:X\to Y$ be a finite morphism (or even, by passage to the limit, an
integral morphism) and $\sF$ an abelian sheaf on $X$. Then $\eR^q f_* \sF=0$,
for all $q>0$.
\end{proposition}

Indeed let $\bar y$ be a geometric point of $Y$, $\widetilde Y$ the
spectrum of the strict localisation of $Y$ at $y$ and
$\widetilde X=X\times_Y \widetilde Y$; by what precedes, it suffices to
show that $\h^q(\widetilde X,\sF) = 0$ for all $q>0$. Now $\widetilde X$ is the
spectrum of a product of strictly henselian local rings (cf.
\cite[I]{ra70}), the functor $\Gamma(\widetilde X,-)$ is exact since every
surjective \'etale $\widetilde X$-scheme admits a section, whence the assertion. 










\subsection{Galois cohomology}\label{I:2-4}

For $X=\spec(K)$ the spectrum of a field, we shall see that \'etale
cohomology is identified with Galois cohomology. 





\subsubsection{}\label{I:2-4-1}

Let us begin with a topological analogy. If $K$ is the function field
of an integral affine algebraic variety $Y=\spec(A)$ over $\dC$, we have
$K=\varinjlim_{f\in A} A[1/f]$.

In other words $X=\varprojlim U$, $U$ running over the set of open subsets of $Y$.
It is known that there exist arbitrarily small Zariski open subsets which for the
classical topology are $K(\pi,1)$'s. We shall therefore not be surprised if
$\spec(K)$ itself is regarded as a $K(\pi,1)$, $\pi$ being the
fundamental group (in the algebraic sense) of $X$, i.e.\ the Galois group of
$\bar K/K$, where $\bar K$ is a separable closure of $K$. 





\subsubsection{}\label{I:2-4-2}

More precisely let $K$ be a field, $\bar K$ a separable closure of $K$ and
$G = \gal(\bar K/K)$ the topological Galois group. To every finite \'etale
$K$-algebra $A$ (finite product of separable extensions of $K$), associate the
finite set $\hom_K(A,\bar K)$. The Galois group $G$ acts on this set through
a discrete (hence finite) quotient. If $A=K[T]/(F)$, it is identified with
the set of roots in $\bar K$ of the polynomial $F$. Galois theory,
in the form given to it by Grothendieck, says that:





\begin{proposition}\label{I:2-4-3}
The functor
\[
  \left\{\begin{array}{c}
          \textnormal{finite \'etale} \\
          \textnormal{$K$-algebras}
        \end{array}\right\}
  \to
  \left\{\begin{array}{c}
          \textnormal{finite sets on which} \\
          \textnormal{$G$ acts continuously}
        \end{array}\right\}
\]
which to an \'etale algebra $A$ associates $\hom_K(A,\bar K)$ is an
anti-equivalence of categories.
\end{proposition}

We deduce an analogous description of sheaves for the \'etale topology on
$\spec(K)$:





\begin{proposition}\label{I:2-4-4}
The functor
\[
  \left\{\begin{array}{c}
          \textnormal{\'Etale sheaves} \\
          \textnormal{on $\spec(K)$}
        \end{array}\right\}
  \to
  \left\{\begin{array}{c}
          \textnormal{sets on which} \\
          \textnormal{$G$ acts continuously}
        \end{array}\right\}
\]
which to a
sheaf $\sF$ associates its fibre $\sF_{\bar K}$ at the geometric point
$\spec(\bar K)$ is an equivalence of categories.
\end{proposition}

We say that $G$ acts continuously on a set $E$ if the stabiliser of every
element of $E$ is an open subgroup of $G$. The functor in the
reverse direction is described in the evident way: let $A$ be a
finite \'etale $K$-algebra, $U=\spec(A)$ and $U(\bar K)=\hom_K(A,\bar K)$ the
corresponding $G$-set; then we have
$\sF(U)= \hom_{G\text{-sets}}(U(\bar K),\sF_{\bar K})$.

In particular, if $X=\spec(K)$, we have $\sF(X) = \sF_{\bar K}^G$. If we
restrict to abelian sheaves, passing to derived functors we obtain canonical
isomorphisms
\[
  \h^q(\et X,\sF) = \h^q(G,\sF_{\bar K})
\]





\subsubsection{Examples}\label{I:2-4-5}

\begin{enumerate}[\indent a)]
  \item To the constant sheaf $\dZ/n$ corresponds $\dZ/n$ with trivial $G$-action.
  \item To the sheaf of $n$-th roots of unity $\dmu_n$ corresponds
    the group $\dmu_n(\bar K)$ of $n$-th roots of unity in
    $\bar K$, with the natural $G$-action.
  \item To the sheaf $\dG_m$ corresponds the group $\bar K^\times$ with the natural
    $G$-action.
\end{enumerate}




















\section{Cohomology of curves}\label{I:3}

In the case of topological spaces, d\'evissages using the K\"unneth formula
and simplicial decompositions allow one to reduce, for
computing cohomology, to the interval $I=[0,1]$ for which we have
$\h^0(I,\dZ)=\dZ$ and $\h^q(I,\dZ)=0$ for $q>0$.

In our case, the d\'evissages will lead to more complicated objects,
namely curves over an algebraically closed field; we shall compute
their cohomology in this chapter. The situation is more complex than in the
topological case since the cohomology groups vanish for $q>2$ only.
The essential ingredient of the computations is the vanishing of the Brauer group of the
function field of such a curve (Tsen's theorem, \ref{I:3-2}). 










\subsection{The Brauer group}\label{I:3-1}

Let us first recall its classical definition:





\begin{definition}\label{I:3-1-1}
Let $K$ be a field and $A$ a finite-dimensional $K$-algebra. We say that $A$
is a central simple algebra over $K$ if the following equivalent
conditions hold:
\begin{enumerate}[\indent a)]
  \item $A$ has no non-trivial two-sided ideal and its centre is $K$.
  \item There exists a finite Galois extension $K'/K$ such that
    $A_{K'} = A\otimes_K K'$ is isomorphic to a square matrix
    algebra over $K'$.
  \item $A$ is $K$-isomorphic to a square matrix algebra over a skew
    field with centre $K$.
\end{enumerate}
\end{definition}

Two such algebras are said to be equivalent if the skew fields
associated to them by c) are $K$-isomorphic. If these algebras have the
same dimension, this amounts to saying that they are $K$-isomorphic. The
tensor product defines, by passage to the quotient, an abelian group
structure on the set of equivalence classes. It is this group that is
classically called \emph{the Brauer group} of $K$ and denoted
$\br(K)$. 





\subsubsection{}\label{I:3-1-2}

We denote by $\br(n,K)$ the set of $K$-isomorphism classes of
$K$-algebras $A$ such that there exists a finite Galois extension $K'$ of
$K$ for which $A_{K'}$ is isomorphic to the algebra $M_n(K')$ of
square $n\times n$ matrices over $K'$. By definition $\br(K)$ is the union of the
subsets $\br(n,K)$ for $n\in\dN$. Let $\bar K$ be an algebraic
closure of $K$ and $G=\gal(\bar K/K)$. The set $\br(n,K)$ is
the set of ``forms'' of $M_n(\bar K)$, hence canonically
isomorphic to $\h^1\left(G,\aut(M_n(\bar K))\right)$.

It is known that every automorphism of $M_n(\bar K)$ is inner. Hence
the group $\aut(M_n(\bar K))$ is identified with the projective linear
group $\pgl(n,\bar K)$ and we have a canonical bijection:
\[
  \theta_n : \br(n,K) \iso \h^1\left(G,\pgl(n,\bar K)\right)\text{.}
\]

On the other hand the exact sequence:
\begin{equation*}\tag{3.1.2.1}\label{I:eq:3-1-2-1}
  1 \to \bar K^\times \to \gl(n,\bar K) \to \pgl(n,\bar K) \to 1\text{,}
\end{equation*}
allows one to define a coboundary operator:
\[
  \Delta_n : \h^1\left(G,\pgl(n,\bar K)\right) \to \h^2(G,\bar K^\times)\text{.}
\]

Composing $\theta_n$ and $\Delta_n$, we obtain a map:
\[
  \delta_n : \br(n,K) \to \h^2(G,\bar K^\times)\text{.}
\]
One checks easily that the maps $\delta_n$ are compatible with each
other and define a group homomorphism:
\[
  \delta : \br(K) \to \h^2(G,\bar K^\times)\text{.}
\]





\begin{proposition}\label{I:3-1-3}
The homomorphism $\delta:\br(K)\to \h^2(G,\bar K^\times)$ is bijective.
\end{proposition}

This follows from the two following lemmas:





\begin{lemma}\label{I:3-1-4}
The map
$\Delta_n:\h^1\left(G,\pgl(n,\bar K)\right)\to \h^2(G,\bar K^\times)$ is
injective.
\end{lemma}

By \cite{se94}, cor.\ to prop.\ I-44, it suffices to check that each
time one twists the exact sequence (\ref{I:eq:3-1-2-1}) by an element of
$\h^1\left(G,\pgl(n,\bar K)\right)$, the $\h^1$ of the middle group is trivial.
This middle group is the group of $\bar K$-points of the multiplicative group
of a central simple algebra $A$ of rank $n^2$ over $K$. To prove that
$\h^1(G,A_{\bar K}^\times) = 0$, one interprets $A^\times$ as the group
of automorphisms of the free $A$-module $L$ of rank $1$, and $\h^1$ as
the set of ``forms'' of $L$ -- $A$-modules of rank $n^2$ over $K$,
automatically free. 





\begin{lemma}\label{I:3-1-5}
Let $\alpha\in \h^2(G,\bar K^\times)$, $K'$ a finite extension of $K$
contained in $\bar K$, $n=[K':K]$, and $G'=\gal(\bar K/K')$. If the image of
$\alpha$ in $\h^2(G',\bar K^\times)$ is zero, then $\alpha$ belongs
to the image of $\Delta_n$.
\end{lemma}

Let us first note that we have: 
\[
  \h^2(G',\bar K^\times) \simeq \h^2\left(G,(\bar K\otimes_K K')^\times\right)\text{.}
\]
(From a geometric point of view, denoting by $x=\spec(K)$, $x'=\spec(K')$ and
$\pi:x'\to x$ the canonical morphism, we have $\eR^q \pi_*(\dG_{m,x'}) = 0$ for
$q>0$ and hence $\h^q(x',\dG_{m,X'})\simeq \h^q(x,\pi_*\dG_{m,X'})$ for
$q\geqslant 0$).

Moreover the choice of a basis of $K'$ as a vector space over $K$
allows one to define a homomorphism
\[
  (\bar K\otimes_K K')^\times \to \gl(n,\bar K)
\]
which, to an element $x$, associates the endomorphism of
multiplication by $x$ on $\bar K\otimes_K K'$. We then have a commutative
diagram with exact rows: 
\[\xymatrix{
  1 \ar[r] 
    & \bar K^\times \ar[r] \ar@{=}[d]
    & (\bar K\otimes_K K')^\times \ar[r] \ar[d] 
    & (\bar K\otimes_K K')^\times / \bar K^\times \ar[r] \ar[d] 
    & 1 \\
  1 \ar[r] 
    & \bar K^\times \ar[r] 
    & \gl(n,\bar K) \ar[r] 
    & \pgl(n,\bar K) \ar[r] 
    & 1
}\]
The lemma follows from the commutative diagram deduced from it by passing
to cohomology:
\[\xymatrix{
  \h^1\left(G,(\bar K\otimes_K K')^\times/\bar K^\times\right) \ar[r] \ar[d] 
    & \h^2(G,\bar K^\times) \ar[r] \ar@{=}[d]
    & \h^2\left(G,(\bar K\otimes_K K')^\times\right) \\
  \h^1\left(G,\pgl(n,\bar K)\right) \ar[r]^-{\Delta_n} 
    & \h^2(G,\bar K^\times) \text{.}
}\]





\begin{proposition}\label{I:3-1-6}
Let $K$ be a field, $\bar K$ an algebraic closure of $K$ and
$G=\gal(\bar K/K)$. Suppose that, for every finite extension $K'$ of $K$, we
have $\br(K')=0$. Then we have:
\begin{enumerate}[\indent i)]
  \item $\h^q(G,\bar K^\times) = 0$ for all $q>0$.
  \item $\h^q(G,\sF) = 0$ for every torsion $G$-module $\sF$ and every
    $q\geqslant 2$.
\end{enumerate}
\end{proposition}

(For the proof, cf. \cite{se94}).










\subsection{Tsen's theorem}\label{I:3-2}





\begin{definition}\label{I:3-2-1}
We say that a field $K$ is $C_1$ if every non-constant homogeneous
polynomial $f(x_1,\dotsc,x_n)$ of degree $d<n$ has a non-trivial zero.
\end{definition}





\begin{proposition}\label{I:3-2-2}
If a field $K$ is $C_1$, then $\br(K)=0$.
\end{proposition}

This amounts to showing that every skew field $D$ with centre $K$ and finite
over $K$ is equal to $K$. Let $r^2$ be the degree of $D$ over $K$ and
$\operatorname{Nrd}:D\to K$ the reduced norm.

(Locally for the \'etale topology over $K$, $D$ is isomorphic -- non
canonically -- to a matrix algebra $M_r$ and the reduced norm
coincides with the determinant map. The latter is well defined,
independently of the chosen isomorphism between $D$ and $M_r$, since every
automorphism of $M_r$ is inner and two similar matrices have the
same determinant. This map, defined locally for the
\'etale topology, descends, by its local uniqueness, to a
map $\operatorname{Nrd}:D\to K$).

The only zero of $\operatorname{Nrd}$ is the zero element of $D$, for if
$x\ne 0$, we have $\operatorname{Nrd}(x)\cdot \operatorname{Nrd}(x^{-1}) = 1$. 
On the other hand, if $\{e_1,\dotsc,e_{r^2}\}$ is a basis of $D$ over $K$ and
$x=\sum x_i e_i$, the function $\operatorname{Nrd}(x)$ is expressed as a homogeneous
polynomial $\operatorname{Nrd}(x_1,\dotsc,x_{r^2})$ of degree $r$ (this is
clear locally for the \'etale topology). Since $K$ is $C_1$, we have
$r^2\leqslant r$, i.e.\ $r=1$ and $D=K$. 





\begin{theorem}[Tsen]\label{I:3-2-3}
Let $k$ be an algebraically closed field and $K$ an extension of
transcendence degree $1$ of $k$. Then $K$ is $C_1$.
\end{theorem}

Suppose first that $K=k(X)$. Let
\[
  f(T) = \sum a_{i_1,\dotsc i_n} T_1^{i_1} \dotsm T_n^{i_n}
\]
be a homogeneous polynomial of degree $d<n$ with coefficients in $k(X)$.
Multiplying the coefficients by a common denominator we may assume
they lie in $k[X]$. Let $\delta=\sup\deg(a_{i_1\dotsc i_n})$. We
seek a non-trivial zero in $k[X]$ by the method of indeterminate
coefficients, writing each $T_i$ ($i=1,\dotsc,n$) as a
polynomial of degree $N$ in $X$. Then the equation $f(T)=0$ becomes a system
of homogeneous equations in the $n\times (N+1)$ coefficients of the polynomials
$T_i(X)$ expressing the vanishing of the coefficients of the polynomial in $X$ obtained by
replacing $T_i$ by $T_i(X)$. This polynomial is of degree $\delta+N D$ at most,
so there are $\delta+N d+1$ equations in $n\times (N+1)$ variables. Since $k$
is algebraically closed this system has a non-trivial solution if
$n(N+1)>N d+\delta+1$, which will be the case for $N$ large enough if $d<n$.

It is clear that, to prove the theorem in the general case, it
suffices to prove it when $K$ is a finite extension of a purely
transcendental extension $k(X)$ of $k$. Let $f(T)=f(T_1,\dotsc,T_n)$ be a
homogeneous polynomial of degree $d<n$ with coefficients in $K$. Let
$a=[K:k(X)]$ and $e_1,\dotsc,e_s$ a basis of $K$ over $k(X)$. Introduce
new variables $U_{i j}$, $s n$ in number, such that
$T_i=\sum U_{i j}e_j$. For the polynomial $f(T)$ to have a non-trivial
zero in $K$, it suffices that the polynomial
$g(X_{i j}) = N_{K/k}(f( T))$ have a non-trivial zero in $k(X)$.
Now $g$ is a polynomial of degree $s d$ in $s n$ variables, whence the
result. 





\begin{corollary}\label{I:3-2-4}
  Let $k$ be an algebraically closed field and $K$ an extension of
transcendence degree $1$ of $k$. Then the \'etale cohomology groups
$\h^q(\spec(K),\dG_m)$ vanish for all $q>0$.
\end{corollary}










\subsection{Cohomology of smooth curves}

Henceforth, and unless expressly stated otherwise, the cohomology
groups under consideration are \'etale cohomology groups. 





\begin{proposition}\label{I:3-3-1}
Let $k$ be an algebraically closed field and $X$ a connected non-singular
projective curve over $k$. Then we have:
\begin{align*}
  \h^0(X,\dG_m) &= k^\times \text{,} \\
  \h^1(X,\dG_m) &= \pic(X) \text{,} \\
  \h^q(X,\dG_m) &= 0 \text{ for $q\geqslant 2$.}
\end{align*}
\end{proposition}

Let $\eta$ be the generic point of $X$, $j:\eta\to X$ the canonical
morphism and $\dG_{m,X}$ the multiplicative group of the function field
$K(X)$. For every closed point $x$ of $X$, let $i_x:x\to X$ be the canonical
immersion and $\dZ_x$ the constant sheaf with value $\dZ$ on $x$. Thus
$j_*\dG_{m,\eta}$ is the sheaf of non-zero meromorphic functions on
$X$ and $\oplus_{x\in X} i_{x_I}\dZ_x$ the sheaf of divisors, so we have an
exact sequence of sheaves:
\begin{equation*}\tag{3.3.1.1}\label{I:eq:3-3-1-1}
\xymatrix{
  0 \ar[r] 
    & \dG_m \ar[r]
    & j_* \dG_{m,\eta} \ar[r]^-{\dv} \ar[r] 
    & \displaystyle\bigoplus_{x\in X} i_{x*} \dZ_x \ar[r] 
    & 0\text{.}
}
\end{equation*}





\begin{lemma}\label{I:3-3-2}
We have $\eR^q j_* \dG_{m,\eta} = 0$ for all $q>0$.
\end{lemma}

It suffices to show that the fibre of this sheaf at every closed point $x$
of $X$ is zero. If $\tilde\sO_{X,x}$ is the henselisation of $X$ at $x$ and $K$
the fraction field of $\tilde\sO_{X,x}$, we have
\[
  \spec(K) = \eta\times_X \spec(\tilde\sO_{X,x})\text{,}
\]
hence $(\eR^qj_*\dG_{m,\eta})_x = \h^q(\spec(K),\dG_m)$.

Now $K$ is an algebraic extension of $k(X)$, hence an extension of
transcendence degree $1$ of $k$: the lemma follows from (\ref{I:3-2-4}). 





\begin{lemma}\label{I:3-3-3}
We have $\h^q(X, j_*\dG_{m,\eta}) = 0$ for all $q>0$.
\end{lemma}

Indeed from (\ref{I:3-3-2}) and the Leray spectral sequence for $j$, we
deduce:
\[
  \h^q(X, j_* \dG_{m,\eta}) = \h^q(\eta,\dG_{m,\eta})
\]
for all $q\geqslant 0$ and the right-hand side is zero for $q>0$ by
(\ref{I:3-2-4}). 





\begin{lemma}\label{I:3-3-4}
We have $\h^q\left(X,\bigoplus_{x\in X} i_{x*} \dZ_x\right) = 0$ for all $q>0$.
\end{lemma}

Indeed for every closed point of $X$, we have $\eR^q i_{x*}\dZ_x = 0$ for $q>0$,
since $i_x$ is a finite morphism (\ref{I:2-3-6}), and
\[
  \h^q(X,i_{x*}\dZ_x) = \h^q(x,\dZ_x)\text{.}
\]
The right-hand side is zero for all $q>0$, since $x$ is the spectrum of an
algebraically closed field (one sees that the lemma holds more
generally for every ``skyscraper'' sheaf on $X$).

We deduce from the preceding lemmas and the exact sequence (\ref{I:eq:3-3-1-1})
the equalities:
\[
  \h^q(X,\dG_m) = 0 \text{ for $q\geqslant 2$,}
\]
and a low-degree cohomology exact sequence:
\[
  1 \to \h^0(X,\dG_m) \to \h^0(X, j_*\dG_{m,\eta}) \to \h^0\left(X,\bigoplus_{x\in X} i_{x*}\dZ_x\right) \to \h^1(X,\dG_m) \to 1
\]
which is none other than the exact sequence:
\[
  1 \to k^\times \to k(X)^\times\to \Div(X) \to \pic(X) \to 1\text{.}
\]
From Proposition \ref{I:3-3-1} we deduce that the cohomology groups of
$X$ with values in $\dZ/n$, $n$ prime to the characteristic of $k$, take
a reasonable value:





\begin{corollary}\label{I:3-3-5}
If $X$ is of genus $g$ and $n$ is invertible in $k$, the
$\h^q(X,\dZ/n)$ vanish for $q>2$, and are free over $\dZ/n$ of rank $1,2g,1$
for $q=0,1,2$. Replacing $\dZ/n$ by the isomorphic group $\dmu_n$, we have
canonical isomorphisms
\begin{align*}
  \h^0(X,\dmu_n) &= \dmu_n \\
  \h^1(X,\dmu_n) &= \pic^0(X)_n \\
  \h^2(X,\dmu_n) &= \dZ/n \text{.}
\end{align*}
\end{corollary}

Since the field $k$ is algebraically closed, $\dZ/n$ is (non-canonically)
isomorphic to $\mu_n$. From the Kummer exact sequence:
\[
  0 \to \mu_n \to \dG_m \to \dG_m \to 0\text{,}
\]
and Proposition \ref{I:3-3-1}, we deduce the equalities:
\[
  \h^q(X,\dZ/n) = 0 \text{ for $q>2$,}
\]
and, in low degree, exact sequences:
\[\xymatrix{
  0 \ar[r] 
    & \h^0(X,\dmu_n) \ar[r] 
    & k^\times \ar[r]^-n 
    & k^\times \ar[r] 
    & 0
}\]
\[\xymatrix{
  0 \ar[r] 
    & \h^1(X,\dmu_n) \ar[r] 
    & \pic(X) \ar[r]^-n 
    & \h^2(X,\dmu_n) \ar[r] 
    & 0\text{.}
}\]
Moreover we have an exact sequence:
\[\xymatrix{
  0 \ar[r] 
    & \pic^0(X) \ar[r] 
    & \pic(X) \ar[r]^-{\deg}
    & \dZ \ar[r] 
    & 0 \text{,}
}\]
and $\pic^0(X)$ is identified with the group of rational points over $k$ of an
abelian variety of dimension $g$, the Jacobian of $X$. In such a
group, multiplication by $n$ is surjective and its kernel is a
free $\dZ/n\dZ$-module of rank $2g$ (since $n$ is invertible in $k$); whence
the corollary.

A clever d\'evissage, using the ``trace method,'' yields as a corollary 





\begin{proposition}[{\cite[IX 5.7]{sga4}}]\label{I:3-3-6}
Let $k$ be an algebraically closed field, $X$ an algebraic curve over $k$
and $\sF$ a torsion sheaf on $X$. Then:
\begin{enumerate}[\indent i)]
  \item We have $\h^q(X,\sF) = 0$ for $q>2$.
  \item If $X$ is affine, we even have $\h^q(X,\sF) = 0$ for $q>1$.
\end{enumerate}
\end{proposition}

For the proof, as well as for the exposition of the ``trace
method,'' we refer to \cite[IX 5]{sga4}.










\subsection{D\'evissages}\label{I:3-4}

To compute the cohomology of varieties of dimension $>1$ one uses
fibrations by curves, which allow one to reduce to studying
morphisms whose fibres are of dimension $\leqslant 1$. This principle
has several variants; let us indicate a few.





\subsubsection{}\label{I:3-4-1}

Let $A$ be a finitely generated $k$-algebra and $a_1,\dotsc,a_n$ generators of
$A$. Putting $X_0=\spec(k)$, $X_i=\spec(k[a_1,\dotsc,a_i])$,
$X_n=\spec(A)$, the canonical inclusions
$k[a_1,\dotsc,a_i]\to k[a_1,\dotsc,a_i,a_{i+1}]$ define morphisms
$X_n \to X_{n-1} \to \cdots \to X_1 \to X_0$ whose fibres are of dimension
$\leqslant 1$. 





\subsubsection{}\label{I:3-4-2}

In the case of a smooth morphism, one can be more precise. An
\emph{elementary fibration} is a morphism of schemes $f:X\to S$ which can
be embedded in a commutative diagram
\[\xymatrix{
  X \ar@{^{(}->}[r]^-j \ar[dr]_-f
    & \bar X \ar[d]^-{\bar f}
    & \ar@{_{(}->}[l]_-i \ar[dl]^-g Y \\
  & S
}\]
satisfying the following conditions:
\begin{enumerate}[\indent i)]
  \item $j$ is a dense open immersion in each fibre and
    $X=\bar X \setminus Y$.
  \item $\bar f$ is smooth and projective, with irreducible geometric fibres
    of dimension $1$.
  \item $g$ is an \'etale covering and no fibre of $g$ is empty.
\end{enumerate}

A \emph{good neighbourhood} relative to $S$ is an $S$-scheme $X$ for which there
exist $S$-schemes $X=X_n,\dotsc,X_0=S$ and elementary fibrations
$f_i:X_i\to X_{i-1}$, $i=1,\dotsc,n$. One can show \cite[XI 3.3]{sga4} that if
$X$ is a smooth scheme over an algebraically closed field $k$, every
rational point of $X$ has an open neighbourhood which is a good neighbourhood
(relative to $\spec(K)$). 





\subsubsection{}\label{I:3-4-3}

One can d\'evisser a proper morphism $f:X\to S$ as follows.
By Chow's lemma, there exists a commutative diagram
\[\xymatrix{
  X \ar[dr]_-f
    & \ar[l]_-\pi \ar[d]^-{\bar f} \bar X \\
  & S
}\]
where $\pi$ and $\bar f$ are projective morphisms, $\pi$ being moreover an
isomorphism over a dense open subset of $X$. Locally over $S$, $\bar X$
is a closed subscheme of a projective space of type $\dP_S^n$.

One d\'evisses the latter by considering the projection
$\varphi:\dP_S^n\to\dP_S^1$ sending the point of homogeneous coordinates
$(x_0,x_1,\dotsc,x_n)$ to $(x_0,x_1)$. It is a rational map
defined outside the closed subset $Y\simeq \dP_S^{n-2}$ of $\dP_S^n$ with
homogeneous equations $x_0=x_1=0$. Let $u:P\to \dP_S^n$ be the blow-up with
centre $Y$; the fibres of $u$ are of dimension $\leqslant 1$. Moreover there exists a
natural morphism $v:P\to \dP_S^1$ extending the rational map
$\varphi$ and $v$ makes $P$ a $\dP_S^1$-scheme locally isomorphic to
projective space of type $\dP^{n-1}$, which in turn can be projected onto a
$\dP^1$, etc. 





\subsubsection{}\label{I:3-4-4}

One can sweep a smooth projective variety $X$ by a \emph{Lefschetz pencil}.
The blow-up $\widetilde X$ of the intersection of the axis of the pencil
with $X$ projects onto $\dP^1$ and the fibres of this projection are the
hyperplane sections of $X$ by the hyperplanes of the pencil.




















\section{Base change theorem for a proper morphism}\label{I:4}










\subsection{Introduction}\label{I:4-1}

This chapter is devoted to the proof and the applications of the 





\begin{theorem}\label{I:4-1-1}
Let $f:X\to S$ be a proper morphism of schemes and $\sF$ a torsion abelian
sheaf on $X$. Then, for every $q\geqslant 0$, the fibre of
$\eR^q f_* \sF$ at a geometric point $s$ of $S$ is isomorphic to the
cohomology $\h^q(X_s,\sF)$ of the fibre $X_s=X\otimes_S \spec k(s)$ of $f$ at
$s$.
\end{theorem}

For $f:X\to S$ a proper continuous and separated map (separated
meaning that the diagonal of $X\times_S X$ is closed) between
topological spaces, and $\sF$ an abelian sheaf on $X$, the analogous result is
well known, and elementary: since $f$ is closed, the $f^{-1}(V)$ for
$V$ a neighbourhood of $s$ form a fundamental system of neighbourhoods of $X_s$, and
one checks that $\h^\bullet(X_s,\sF) = \varinjlim_U \h^\bullet(U,\sF)$, for $U$ running over
neighbourhoods of $X_s$. In practice, $X_s$ even has a fundamental system
$\cU$ of neighbourhoods $U$ of which it is a deformation retract and, for $\sF$
constant, one thus has $\h^\bullet(X_s,\sF)=\h^\bullet(U,\sF)$. In picturesque terms: the
special fibre swallows the general fibre.

In the case of schemes the proof is more delicate and it is
indispensable to assume that $\sF$ is torsion (\cite[XII.2]{sga4}). In view
of the description of the fibres of $\eR^q f_* \sF$ (\ref{I:2-3-3}), Theorem
(\ref{I:4-1-1}) is essentially equivalent to the 

 
  
   
    
\begin{theorem}\label{I:4-1-2}
Let $A$ be a strictly henselian local ring and $S=\spec(A)$. Let
$f:X\to S$ be a proper morphism and $X_0$ the closed fibre of $f$. Then, for
every torsion abelian sheaf $\sF$ on $X$ and every $q\geqslant 0$, we
have $\h^q(X,\sF) \iso \h^q(X_0,\sF)$.
\end{theorem}

By passage to the limit one sees that it suffices to prove the theorem
when $A$ is the strict henselisation of a $\dZ$-algebra of finite type at
a prime ideal. One treats first the case $q=0$ or $1$ and $\sF=\dZ/n$
(\ref{I:4-2}). An argument based on the notion of constructible sheaf
(\ref{I:4-3}) shows moreover that it suffices to consider the case where $\sF$ is
constant. On the other hand the d\'evissage (\ref{I:3-4-3}) allows one to assume that
$X_0$ is a curve; in this case it only remains to prove the
theorem for $q=2$ (\ref{I:4-4}).

Among other applications (\ref{I:4-6}), the theorem makes it possible to define the
notion of cohomology with proper support (\ref{I:4-5}). 










\subsection{Proof for \texorpdfstring{$q=0$}{q=1} or \texorpdfstring{$1$}{1} and \texorpdfstring{$\sF=\dZ/n$}{\sF=Z/n}}\label{I:4-2}

The result for $q=0$ and $\sF$ constant is equivalent to the following
proposition (Zariski's connectedness theorem):





\begin{proposition}\label{I:4-2-1}
Let $A$ be a noetherian henselian local ring and $S=\spec(A)$.
Let $f:X\to S$ be a proper morphism and $X_0$ the closed fibre of $f$. Then
the sets of connected components $\pi_0(X)$ and $\pi_0(X_0)$ are in
bijection.
\end{proposition}

It amounts to the same to show that the sets of simultaneously open
and closed subsets $\operatorname{Of}(X)$ and $\operatorname{Of}(X_0)$ are
in bijection. It is known that the set $\operatorname{Of}(X)$ is in
bijection with the set of idempotents of $\Gamma(X,\sO_X)$, and likewise
$\operatorname{Of}(X_0)$ is in bijection with the set of
idempotents of $\Gamma(X_0,\sO_{X_0})$. It is therefore a matter of showing that
the canonical map
\[
  \operatorname{Idem} \Gamma(X,\sO_X) \to \operatorname{Idem}\Gamma(X_0,\sO_{X_0})
\]
is bijective.

We denote by $\fm$ the maximal ideal of $A$, by $\Gamma(X,\sO_X)^\wedge$ the
completion of $\Gamma(X,\sO_X)$ for the $\fm$-adic topology and, for every
integer $n\geqslant 0$, $X_n=X\otimes_A A/\fm^{n+1}$. By the finiteness theorem
for proper morphisms \cite[III.3.2]{ega}, $\Gamma(X,\sO_X)$ is a
finite $A$-algebra; since $A$ is henselian, it follows that
the canonical map
\[
  \operatorname{Idem} \Gamma(X,\sO_X) \to \operatorname{Idem}\Gamma(X,\sO_X)^\wedge
\]
is bijective. In particular the canonical map
\[
  \operatorname{Idem} \Gamma(X,\sO_X)^\wedge \to \varprojlim \operatorname{Idem}\Gamma(X_n,\sO_{X_n})
\]
is bijective. But, since $X_n$ and $X_0$ have the same underlying
topological space, the canonical map
\[
  \operatorname{Idem}\Gamma(X_n,\sO_{X_n}) \to \operatorname{Idem} \Gamma(X_0,\sO_{X_0})
\]
is bijective for all $n$, which completes the proof.

Since $\h^1(X,\dZ/n)$ is in bijection with the set of isomorphism classes
of Galois \'etale coverings of $X$ with group $\dZ/n$,
the theorem for $q=1$ and $\sF=\dZ/n$ follows from the next proposition. 





\begin{proposition}\label{I:4-2-2}
Let $A$ be a noetherian henselian local ring and $S=\spec(A)$. Let
$f:X\to S$ be a proper morphism and $X_0$ the closed fibre of $f$. Then the
restriction functor
\[
  \operatorname{Rev.et}(X) \to \operatorname{Rev.et}(X_0)
\]
is an equivalence of categories.
\end{proposition}

(If $X_0$ is connected and a geometric point of $X_0$ has been chosen as
base point, this amounts to saying that the canonical map
$\pi_1(X_0)\to \pi_1(X)$ on (profinite) fundamental groups is
bijective).

Proposition (\ref{I:4-2-1}) shows that this functor is fully faithful. Indeed,
if $X'$ and $X''$ are two \'etale coverings of $X$, an
$X$-morphism from $X'$ to $X''$ is determined by its graph, which is an
open and closed subset of $X'\times_X X''$.

It is therefore a matter of showing that every \'etale covering $X_0'$ of $X_0$
extends to an \'etale covering of $X$. It is known that \'etale coverings do
not depend on nilpotents \cite[ch.1]{sga1}, hence
$X_0'$ lifts uniquely to an \'etale covering $X_n'$ of
$X_n$ for all $n\geqslant 0$, i.e.\ to an \'etale covering
$\fX'$ of the formal scheme $\fX$, the completion of $X$ along $X_0$. By
Grothendieck's algebraisation theorem for formal coherent sheaves
(existence theorem, \cite[III.5]{ega}), $\fX'$ is the
formal completion of an \'etale covering $\bar X'$ of
$\bar X=X\otimes_A \hat A$.

By passage to the limit, it suffices to prove the proposition in the case
where $A$ is the henselisation of a $\dZ$-algebra of finite type. One can then
apply Artin's approximation theorem to the functor
$\sF:(\textnormal{$A$-algebras}) \to (\textnormal{sets})$ which, to an
$A$-algebra $B$, associates the set of isomorphism classes of
\'etale coverings of $X\otimes_A B$. Indeed this functor is locally of
finite presentation: if $B_i$ is a filtered inductive system of
$A$-algebras and $B=\varinjlim B_i$, we have $\sF(B) = \varinjlim \sF(B_i)$.
By Artin's theorem, given an element
$\xi\in \sF(\hat A)$, in this case the isomorphism class of $\bar X'$, there
exists $\xi\in \sF(A)$ having the same image as $\bar\xi$ in $\sF(A/\fm)$.
In other words there exists an \'etale covering $X'$ of $X$ whose restriction
to $X_0$ is isomorphic to $X_0'$. 










\subsection{Constructible sheaves}\label{I:4-3}

In this paragraph, we consider a \emph{noetherian} scheme $X$ and by a
sheaf on $X$ we mean an \emph{abelian} sheaf on $\et X$. 





\begin{definition}\label{I:4-3-1}
We say that a sheaf $\sF$ on $X$ is \emph{locally constant constructible}
(abbreviated l.c.c.) if it is represented by an \'etale covering of
$X$.
\end{definition}





\begin{definition}\label{I:4-3-2}
We say that a sheaf $\sF$ on $X$ is \emph{constructible} if it satisfies the
following equivalent conditions:
\begin{enumerate}[\indent (i)]
  \item There exists a finite surjective family of subschemes $X_i$ of $X$
     such that the restriction of $\sF$ to $X_i$ is l.c.c..
  \item There exists a finite family of finite morphisms $p_i:X_i'\to X$, for
    each $i$ a finite constant sheaf ( = defined by a finite
    abelian group) $C_i$ on $X_i'$, and a monomorphism
    $\sF\to \prod p_{i*} C_i$.
\end{enumerate}
\end{definition}

One checks easily that the category of constructible sheaves on $X$
is an abelian category. Moreover, if $u:\sF\to \sG$ is a homomorphism of
sheaves and $\sF$ is constructible, the sheaf $\im(u)$ is constructible.





\begin{lemma}\label{I:4-3-3}
Every torsion sheaf $\sF$ is a filtered inductive limit of constructible
sheaves.
\end{lemma}

Indeed, if $j:U\to X$ is an \'etale scheme of finite type over $X$, an
element $\xi\in \sF(U)$ with $n\xi=0$ defines a homomorphism of
sheaves $j_! \const{\dZ/n}\to \sF$ whose image (the smallest subsheaf of
$\sF$ of which $\xi$ is a local section) is a constructible subsheaf of $\sF$.
It is clear that $\sF$ is the inductive limit of such subsheaves.





\begin{definition}\label{I:4-3-4}
Let $\cC$ be an abelian category and $T$ a functor defined on $\cC$
with values in the category of abelian groups. We say that $T$ is
\emph{effaceable} in $\cC$ if, for every object $A$ of $\cC$ and every
$\alpha\in T(A)$, there exists a monomorphism $u:A\to M$ in $\cC$ such that
$T(u)\alpha = 0$.
\end{definition}





\begin{lemma}\label{I:4-3-5}
The functors $\h^q(X,-)$ for $q>0$ are effaceable in the category of
constructible sheaves on $X$.
\end{lemma}

It suffices to note that, if $\sF$ is a constructible sheaf, there necessarily
exists an integer $n>0$ such that $\sF$ is a sheaf of
$\dZ/n$-modules. Then there exists a monomorphism $\sF\hookrightarrow \sG$, where
$\sG$ is a sheaf of $\dZ/n$-modules with $\h^q(X,\sG) = 0$ for all $q>0$.
One can for example take for $\sG$ the Godement resolution
$\prod_{x\in X} i_{x*} \sF_{\bar x}$, where $x$ runs over the points of $X$ and
$i_x:\bar x\to X$ is a geometric point centred at $X$. By
(\ref{I:4-3-3}) $\sG$ is an inductive limit of constructible sheaves, whence the
lemma, since the functors $\h^q(X,-)$ commute with inductive limits. 





\begin{lemma}\label{I:4-3-6}
Let $\varphi^\bullet:T^\bullet\to {T'}^\bullet$ be a morphism of cohomological
functors defined on an abelian category $\cC$ with values
in the category of abelian groups. Suppose that $T^q$ is effaceable
for $q>0$ and let $\cE$ be a set of objects of $\cC$ such that every object
of $\cC$ is contained in an object belonging to $\cE$. Then the
following conditions are equivalent:
\begin{enumerate}[\indent (i)]
  \item $\varphi^q(A)$ is bijective for all $q\geqslant 0$ and all
    $A\in\ob \cC$.
  \item $\varphi^0(M)$ is bijective and $\varphi^q(M)$ surjective for all
    $q>0$ and all $M\in\cE$.
  \item $\varphi^0(A)$ is bijective for all $A\in\ob \cC$ and ${T'}^q$ is
    effaceable for all $q>0$.
\end{enumerate}
\end{lemma}

The proof is by induction on $q$ and presents no difficulty. 





\begin{proposition}\label{I:4-3-7}
Let $X_0$ be a subscheme of $X$. Suppose that, for every $n\geqslant 0$ and
for every scheme $X'$ finite over $X$, the canonical map
\[
  \h^q(X',\dZ/n) \to \h^q(X_0',\dZ/n)\text{,}
\]
where $X_0' = X'\times_X X_0$, is bijective for $q=0$ and surjective for
$q>0$. Then, for every torsion sheaf $\sF$ on $X$ and every
$q\geqslant 0$, the canonical map
\[
  \h^q(X,\sF) \to \h^q(X_0,\sF)
\]
is bijective.
\end{proposition}

By passage to the limit, it suffices to prove the assertion for $\sF$
constructible. One applies Lemma (\ref{I:4-3-6}) taking for $\cC$ the
category of constructible sheaves on $X$, $T^q=\h^q(X,-)$,
${T'}^q=\h^q(X_0,-)$, and $\cE$ the set of constructible sheaves of the
form $\prod p_{i*} C_i$, where $p_i:X_i'\to X$ is a finite morphism and $C_i$
a finite constant sheaf on $X_i'$.










\subsection{End of the proof}\label{I:4-4}

By the method of fibration by curves (\ref{I:3-4-3}), we reduce to
proving the theorem in relative dimension $\leqslant 1$. By
the preceding paragraph, it will suffice to show that, if $S$ is the spectrum
of a strictly henselian noetherian local ring, $f:X\to S$ a
proper morphism whose closed fibre $X_0$ is of dimension $\leqslant 1$
and $n$ an integer $\geqslant 0$, the canonical homomorphism
\[
  \h^q(X,\dZ/n) \to \h^q(X_0,\dZ/n)
\]
is bijective for $q=0$ and surjective for $q>0$.

The cases $q=0$ and $1$ have been seen above and we have $\h^q(X_0,\dZ/n) = 0$
for $q\geqslant 3$; hence it suffices to treat the case $q=2$. One may
of course assume that $n$ is a prime power. If
$n=p^r$, where $p$ is the characteristic of the residue field of $S$,
Artin-Schreier theory shows that one has $\h^2(X_0,\dZ/p^r) = 0$. If
$n=\ell^r$, $\ell\ne p$, one deduces from Kummer theory a commutative
diagram
\[\xymatrix{
  \pic(X) \ar[d] \ar[r]^-\alpha
    & \h^2(X,\dZ/\ell^r) \ar[d] \\
  \pic(X_0) \ar[r]^-\beta
    & \h^2(X_0,\dZ/\ell^r)
}\]
where the map $\beta$ is surjective (we have seen this in Chapter \ref{I:3} for
a smooth curve over an algebraically closed field, but similar arguments
apply to any curve over a separably closed field).

To conclude, it will therefore suffice to show:





\begin{proposition}\label{I:4-4-1}
Let $S$ be the spectrum of a henselian noetherian local ring and $f:X\to S$
a proper morphism whose closed fibre $X_0$ is of dimension
$\leqslant 1$. Then the canonical restriction map
\[
  \pic(X)\to \pic(X_0)
\]
is surjective (it suffices moreover that the morphism $f$ be
\emph{separated} of finite type).
\end{proposition}

To simplify the proof, we assume that $X$ is integral,
although this is not necessary. Every invertible sheaf on $X_0$ is
associated to a Cartier divisor (since $X_0$ is a curve, hence
quasi-projective), so it suffices to show that the canonical map
$\Div(X)\to\Div(X_0)$ is surjective.

Every divisor on $X_0$ is a linear combination of divisors whose
support is concentrated at a single non-isolated closed point of $X_0$. Let
$x$ be such a point, $t_0\in\sO_{X_0,x}$ a regular non-invertible element
of $\sO_{X_0,x}$ and $D_0$ the divisor concentrated at $x$ with local equation
$t_0$. Let $U$ be an open neighbourhood of $x$ in $X$ such that there exists a
section $t\in \Gamma(U,\sO_U)$ lifting $t_0$. Let $Y$ be the closed subset of $U$
with equation $t=0$; shrinking $U$ if necessary, we may assume that $x$ is the
only point of $Y\cap X_0$. Then $Y$ is quasi-finite over $S$ at $x$;
since $S$ is the spectrum of a henselian local ring, it follows that
$Y=Y_1\amalg Y_2$, where $Y_1$ is finite over $S$ and $Y_2$ does not meet
$X_0$. Moreover, since $X$ is separated over $S$, $Y_1$ is closed in $X$.

Shrinking $U$ to a smaller open neighbourhood of $x$, we may
assume that $Y=Y_1$, i.e.\ that $Y$ is closed in $X$. We then
define a divisor $D$ on $X$ lifting $D_0$ by putting $D|X\setminus Y=0$ and
$D|U=\dv(t)$, which makes sense since $t$ is invertible on $U\setminus Y$. 





\subsubsection{Remark}\label{I:4-4-2}

In the case where $f$ is proper, one could also give a proof
of the same style as that of Proposition (\ref{I:4-2-2}). Indeed, since
$X_0$ is a curve, there is no obstruction to lifting an invertible
sheaf on $X_0$ to the infinitesimal neighbourhoods $X_n$ of $X_0$, hence to the
formal completion $\fX$ of $X$ along $X_0$. One concludes by
successively applying Grothendieck's existence theorem and
Artin's approximation theorem. 










\subsection{Cohomology with proper support}\label{I:4-5}





\begin{definition}\label{I:4-5-1}
Let $X$ be a separated scheme of finite type over a field $k$. By a
theorem of Nagata, there exist a scheme $\bar X$ proper over $k$ and an
open immersion $j:X\to\bar X$. For every torsion sheaf $\sF$ on $X$ we
denote by $j_! \sF$ the extension by $0$ of $\sF$ to $\bar X$ and define the
\emph{cohomology groups with proper support}
\[
  \h_c^q(X,\sF)=\h^q(\bar X, j_! \sF) \text{.}
\]
\end{definition}

Let us show that this definition is independent of the chosen
compactification $j:X\to\bar X$. Let $j_1:X\to \bar X_1$ and $j_2:X\to \bar X_2$
be two compactifications. Then $X$ maps to $\bar X_1\times \bar X_2$ by
$x\mapsto (j_1(x),j_2(x))$ and the closed image $\bar X_3$ of $X$ under this
map is a compactification of $X$. We thus have a commutative diagram
\[\xymatrix{
  & \bar X_1 \\
  X \ar[ur]^-{j_1} \ar[rr]^-{(j_1,j_2)} \ar[dr]_-{j_2}
    & & \bar X_3 \ar[ul]_-{p_1} \ar[dl]^-{p_2} \\
  & \bar X_2
}\]
where $p_1$ and $p_2$, the restrictions of the natural projections to
$\bar X_3$, are proper morphisms. It therefore suffices to treat the case where
we have a commutative diagram
\[\xymatrix{
  X \ar[r]^-{j_2} \ar[dr]_-{j_1}
    & \bar X_2 \ar[d]^-p \\
  & \bar X_1
}\]
with $p$ a proper morphism. 





\begin{lemma}\label{I:4-5-2}
We have $p_*(j_{2!} \sF) = j_{1!} \sF$ and $\eR^q p_*(j_{2!} \sF) = 0$, for $q>0$.
\end{lemma}

Let us note at once that the lemma suffices to conclude. Using the
Leray spectral sequence for $p$, we deduce that, for every
$q\geqslant 0$,
\[
  \h^q(\bar X_2, j_{2!} \sF) = \h^q(\bar X_1, j_{1!} \sF) \text{.}
\]

To prove the lemma, we argue fibre by fibre using the base change
theorem (\ref{I:4-1-1}) for $p$. The result is immediate, since,
over a point of $X$, $p$ is an isomorphism and, over a point of
$\bar X_1\setminus X$, $j_{2!} \sF$ is zero on the fibre of $p$. 





\subsubsection{}\label{I:4-5-3}

Similarly, if $f:X\to S$ is a separated morphism of finite type of
noetherian schemes, there exist a proper morphism $\bar f:\bar X\to S$ and an
open immersion $j:X\to \bar X$. One then defines the higher direct images
with proper support $\eR^q f_!$ by putting, for every torsion sheaf
$\sF$ on $X$
\[
  \eR^q f_! \sF = \eR^q f_*(j_! \sF) \text{.}
\]
One checks as before that this definition is independent of
the chosen compactification. 





\begin{theorem}\label{I:4-5-4}
Let $f:X\to S$ be a separated morphism of finite type of noetherian schemes
and $\sF$ a torsion sheaf on $X$. Then the fibre of $\eR^q f_! \sF$ at a
geometric point $s$ of $S$ is isomorphic to the cohomology with proper
support $\h_c^q(X_s,\sF)$ of the fibre $X_s$ of $f$ at $s$.
\end{theorem}

This is a simple variant of the base change theorem for a proper
morphism (\ref{I:4-1-1}). More generally, if
\[\xymatrix{
  X \ar[d]^-f
    & X' \ar[l]_-{g'} \ar[d]^-{f'} \\
  S
    & S' \ar[l]^-g
}\]
is a cartesian diagram, we have a canonical isomorphism
\begin{equation*}\tag{4.5.4.1}\label{I:eq:4-5-4-1}
  g^* (\eR^q f_! \sF) \simeq \eR^q f_!'({g'}^* \sF) \text{.}
\end{equation*}










\subsection{Applications}\label{I:4-6}





\begin{theorem}[vanishing]\label{I:4-6-1}
Let $f:X\to S$ be a separated morphism of finite type whose fibres are of
dimension $\leqslant n$ and $\sF$ a torsion sheaf on $X$. Then we have
$\eR^q f_! \sF = 0$ for $q>2 n$.
\end{theorem}

By the base change theorem, we may assume that $S$ is
the spectrum of a separably closed field. If $\dim X = n$, there exists an
affine open $U$ of $X$ with $\dim(X\setminus U) < n$; we then have an exact
sequence $0\to \sF_U \to \sF \to \sF_{X\setminus U} \to 0$ and, by induction on $n$,
it suffices to prove the theorem for $X=U$ affine. Then the method of
fibration by curves (\ref{I:3-4-1}) and the base change
theorem allow one to reduce to a curve over a separably closed field
for which the desired result follows from Tsen's theorem
(\ref{I:3-3-6}). 





\begin{theorem}[finiteness]\label{I:4-6-2}
Let $f:X\to S$ be a separated morphism of finite type and $\sF$ a
constructible sheaf on $X$. Then the sheaves $\eR^q f_! \sF$ are constructible.
\end{theorem}

We consider only the case where $\sF$ is annihilated by an integer invertible
on $X$.

To prove the theorem one reduces to the case where $\sF$ is the constant sheaf
$\const{\dZ/n}$ and $f:X\to S$ is a proper smooth morphism
whose fibres are geometrically connected curves of genus $g$.
For $n$ invertible on $X$, the sheaves $\eR^q f_* \sF$ are then locally
free of finite rank, zero for $q>2$ (\ref{I:4-6-1}). Replacing $\dZ/n$ by the
locally isomorphic sheaf (over $S$) $\dmu_n$, we have canonically 
\begin{equation*}\tag{4.6.2.1}\label{I:eq:4-6-2-1}
\begin{aligned}
  \eR^0 f_* \dmu_n &= \dmu_n \\
  \eR^1 f_* \dmu_n &= \const\pic(X/S)_n \\
  \eR^2 f_* \dmu_n &= \dZ/n \text{.}
\end{aligned}
\end{equation*}





\begin{theorem}[comparison with classical cohomology]\label{I:4-6-3}
Let $f:X\to S$ be a separated morphism of schemes of finite type over $\dC$,
and $\sF$ a torsion sheaf on $X$. Denote by a superscript $\an{(-)}$ the
functor passing to the usual topological spaces, and by $\eR^q \an f_!$ the
derived functors of the proper direct image functor by $\an f$. We have
\[
  \an{(\eR^q f_! \sF)} \simeq \eR^q \an f_! \an \sF \text{.}
\]
In particular, for $S=$ a point and $\sF$ the constant sheaf $\dZ/n$,
\[
  \h_c^q(X,\dZ/n) \simeq \h_c^q(\an X,\dZ/n)\text{.}
\]
\end{theorem}

D\'evissages using the base change theorem reduce us to the
case where $X$ is a smooth proper curve, $S=$ a point, and
$\sF=\dZ/n$. The cohomology groups under consideration are then zero for
$q\ne 0,1,2$, and one invokes \cite{se55}: indeed, if $X$ is proper over $\dC$,
we have $\pi_0(X) = \pi_0(\an X)$ and $\pi_1(X) = $ the profinite completion of
$\pi_1(\an X)$, whence the assertion for $q=0,1$. For $q=2$, one uses the
Kummer exact sequence and the fact that, by \cite{se55} again,
$\pic(X)=\pic(\an X)$. 





\begin{theorem}[cohomological dimension of affine schemes]\label{I:4-6-4}
Let $X$ be an affine scheme of finite type over a separably closed field and
$\sF$ a torsion sheaf on $X$. Then $\h^1(X,\sF) = 0$ for
$q>\dim(X)$.
\end{theorem}

For the very pretty proof we refer to \cite{sga4}, XIV \S 2 and 3. 





\subsubsection{Remark}\label{I:4-6-5}

This theorem is in some sense a substitute for Morse theory.
Consider indeed the classical case where $X$ is smooth and affine over $\dC$,
embedded in an affine space of type $\dC^N$. Then, for almost every point
$p\in\dC^N$, the ``distance to $p$'' function on $X$ is a Morse function
and the indices of its critical points are smaller than $\dim(X)$. Thus
$X$ is obtained by attaching handles of index smaller than $\dim(X)$, whence
the classical analogue of (\ref{I:4-6-4}). 




















\section{Local acyclicity of smooth morphisms}\label{I:5}

Let $X$ be a complex analytic variety and $f:X\to D$ a morphism from
$X$ to the disc. Denote by $[0,t]$ the closed straight segment with endpoints
$0$ and $t$ in $D$ and by $(0,t]$ the half-open segment. If $f$ is \emph{smooth},
the inclusion
\[
  j : f^{-1}\left((0,t]\right) \hookrightarrow f^{-1}\left([0,t]\right)
\]
is a homotopy equivalence; one can push the special fibre
$X_0=f^{-1}(0)$ into $f^{-1}([0,t])$.

In practice, for $t$ small enough, $f^{-1}((0,t])$ will be a fibre bundle over
$(0,t]$ so that the inclusion
\[
  X_t = f^{-1}(t) \hookrightarrow f^{-1}\left((0,t]\right)
\]
will also be a homotopy equivalence. One then calls \emph{cospecialisation
morphism} the homotopy class of maps:
\[
  \cosp : X_0 \hookrightarrow f^{-1}\left([0,t]\right) \xleftarrow\sim f^{-1}\left((0,t]\right) \xleftarrow\sim X_t \text{.}
\]
This construction can be expressed picturesquely by saying that, for a
smooth morphism, the general fibre swallows the special fibre.

Assume no longer that $f$ is necessarily smooth (but assume that
$f^{-1}((0,t])$ is a fibre bundle over $(0,t]$). One can still define a map
$\cosp^\bullet$ in cohomology as soon as $j_* \dZ=\dZ$ and
$\eR^q j_* \dZ = 0$ for $q>0$. Under these hypotheses, the Leray spectral
sequence for $j$ shows that
\[
  \h^\bullet\left(f^{-1}\left([0,t]\right),\dZ\right) \iso \h^\bullet\left(f^{-1}\left((0,t]\right),\dZ\right)
\]
and $\cosp^\bullet$ is the composite morphism:
\[
  \cosp^\bullet : \h^\bullet(X_t,\dZ)
    \xleftarrow\sim \h^\bullet\left(f^{-1}\left((0,t]\right),\dZ\right)
    \xleftarrow\sim \h^\bullet\left(f^{-1}\left([0,t]\right),\dZ\right)
    \to \h^\bullet(X_0,\dZ)
\]
The fibre of $\eR^q j_* \dZ$ at a point $x\in X_0$ is computed as follows. In
an ambient space take a ball $B_\varepsilon$ of centre $x$ and sufficiently
small radius $\varepsilon$, and for $\eta$ small enough put
$E=X\cap B_\varepsilon \cap f^{-1}(\eta t)$; this is the \emph{variety of
vanishing cycles} at $x$. We have
\[
  (\eR^q j_* \dZ)_x \xleftarrow\sim \h^q\left(X\cap B_\varepsilon\cap f^{-1}\left((0,\eta t]\right),\dZ\right) \xrightarrow\sim \h^q(E,\dZ)
\]
and the cospecialisation morphism is defined in cohomology as soon as the
vanishing cycle varieties are acyclic ($\h^0(E,\dZ)=\dZ$ and
$\h^q(E,\dZ) = 0$ for $q>0$), which is expressed by saying that $f$ is
\emph{locally acyclic}.

This chapter is devoted to the analogue of this situation for a smooth
morphism of schemes in \'etale cohomology. However in this setting it is
indispensable to restrict to torsion coefficients of order \emph{prime to the
residue characteristics}. Paragraph \ref{I:5-1} is devoted
to generalities on locally acyclic morphisms and
cospecialisation maps. In paragraph \ref{I:5-2}, we prove
that a smooth morphism is locally acyclic. In paragraph \ref{I:5-3}, we
combine this result with those of the preceding chapter to deduce two
applications: a theorem on the specialisation of cohomology groups (the
cohomology of the geometric fibres of a proper smooth morphism is
locally constant) and a base change theorem by a smooth morphism.

In all that follows, we fix an integer $n$ and ``scheme'' means ``scheme
on which $n$ is invertible.'' ``Geometric point'' will always mean
``algebraic geometric point'' (\ref{I:2-3-1}) $x:\spec(k)\to X$, with $k$
algebraically closed. 










\subsection{Locally acyclic morphisms}\label{I:5-1}





\subsubsection{Notation}\label{I:5-1-1}

Given a scheme $S$ and a geometric point $s$ of $S$, we
denote by $\widetilde S^s$ the spectrum of the strict localisation of $S$ at $s$. 





\begin{definition}\label{I:5-1-2}
We say that a geometric point $t$ of $S$ is a \emph{generisation} of
$s$ if it is defined by an algebraic closure of the residue field of a
point of $\widetilde S^s$. We also say that $s$ is a \emph{specialisation}
of $t$ and call \emph{specialisation arrow} the $S$-morphism
$t\to \widetilde S^s$.
\end{definition}





\begin{definition}\label{I:5-1-3}
Let $f:X\to S$ be a morphism of schemes. Let $s$ be a geometric point of
$S$, $t$ a generisation of $s$, $x$ a geometric point of $X$
above $s$ and
$\widetilde X_t^x = \widetilde X^x\times_{\widetilde S^s} t$. Then we say that
$\widetilde X_t^x$ is a \emph{vanishing cycle variety} of $f$ at
the point $x$.

We say that $f$ is \emph{locally acyclic} if the reduced cohomology of
every vanishing cycle variety $\widetilde X_t^x$ is zero:
\begin{equation*}\tag{5.1.3.1}\label{I:eq:5-1-3-1}
  \hat\h^\bullet(\widetilde X_t^x,\dZ/n) = 0\text{,}
\end{equation*}
i.e.\ $\h^0\left(X_t^x,\dZ/n\right) = \dZ/n$ and
$\h^q(\widetilde X_t^x,\dZ/n) = 0$ for $q>0$.
\end{definition}





\begin{lemma}\label{I:5-1-4}
Let $f:X\to S$ be a locally acyclic morphism and $g:S'\to S$ a
quasi-finite morphism (or a projective limit of quasi-finite morphisms). Then
the morphism $f':X'\to S'$ deduced from $f$ by base change is locally
acyclic.
\end{lemma}

One checks indeed that every vanishing cycle variety of $f'$ is
a vanishing cycle variety of $f$.	





\begin{lemma}\label{I:5-1-5}
Let $f:X\to S$ be a locally acyclic morphism. For every geometric point
$t$ of $S$, giving rise to a cartesian diagram
\[\xymatrix{
  X_t \ar[r]^-{\varepsilon'} \ar[d]
    & X \ar[d]^-f \\
  t \ar[r]^-\varepsilon
    & S
}\]
we have $\varepsilon_*' \dZ/n = f^*\varepsilon_* \dZ/n$ and $\eR^q \varepsilon_*'\dZ/n=0$
for $q>0$.
\end{lemma}

Let $\bar S$ be the closure of $\varepsilon(t)$, $S'$ the normalisation of $\bar S$
in $k(t)$, and consider the cartesian diagram
\[\xymatrix{
  X_t \ar[r]^-{f'} \ar[d]
    & X' \ar[r]^-{\alpha'} \ar[d]^-{f'}
    & X \ar[d]^-f \\
  t \ar[r]^-i
    & S' \ar[r]^-\alpha
    & S
}\]
The local rings of $S'$ are normal with separably closed fraction field.
They are therefore strictly henselian, and the local acyclicity of $f'$
(\ref{I:5-1-4}) gives $f_*'\dZ/n = \dZ/n$, $\eR^qi_*'\dZ/n = 0$ for $q>0$.
Since $\alpha$ is integral, we then have
\[
  \eR^q\varepsilon_*'\dZ/n
    = \alpha_*' \eR^q i_*' \dZ/n
    = \alpha_*' {f'}^* \eR^q i_* \dZ/n
    = f^* \alpha_* \eR^q i_* \dZ/n
    = f^* \eR^q \varepsilon_* \dZ/n \text{,}
\]
and the lemma. 





\subsubsection{}\label{I:5-1-6}

Given a locally acyclic morphism $f:X\to S$ and a specialisation arrow
$t\to \widetilde S^s$, we shall define canonical
homomorphisms, called cospecialisation arrows
\[
  \cosp^\bullet : \h^\bullet(X_t,\dZ/n) \to \h^\bullet(X_s,\dZ/n)\text{,}
\]
relating the cohomology of the general fibre $X_t=X\times_S t$ and that of the
special fibre $X_s=X\times_X s$.

Consider the cartesian diagram
\[\xymatrix{
  X_t \ar[d] \ar[r]^-{\varepsilon'}
    & \widetilde X \ar[d]^-{f'}
    & \ar[l] X_s \ar[d] \\
  t \ar[r]^-\varepsilon
    & \widetilde S^s
    & \ar[l] s
}\]
deduced from $f$ by base change. By (\ref{I:5-1-4}), $f'$ is still
locally acyclic. From the definition of local acyclicity, one deduces
at once that the restriction to $X_s$ of the sheaves $\eR^q \varepsilon_*' \dZ/n$
is $\dZ/n$ for $q=0$, and $0$ for $q>0$. By (\ref{I:5-1-5}) one even knows that
$\eR^q \varepsilon_*'\dZ/n = 0$ for $q>0$. One defines $\cosp^\bullet$ as the
composite arrow 
\begin{equation*}\tag{5.1.6.1}\label{I:eq:5-1-6-1}
  \h^\bullet(X_t,\dZ/n) \simeq \h^\bullet(X,\varepsilon_*' \dZ/n) \to \h^\bullet(X_s,\dZ/n) \text{.}
\end{equation*}

\paragraph{Variant}
Let $\bar S$ be the closure of $\varepsilon(t)$ in $\widetilde S^s$, $S'$ the
normalisation of $\bar S$ in $k(t)$ and $X'/S'$ deduced from $X/S$ by base
change. The diagram \eqref{I:eq:5-1-6-1} can also be written
\[
  \h^\bullet(X_t,\dZ/n) \simeq \h^\bullet(X',\dZ/n) \to \h^\bullet(X_s,\dZ/n) \text{.}
\]





\begin{theorem}\label{I:5-1-7}
Let $S$ be a locally noetherian scheme, $s$ a geometric point of
$S$ and $f:X\to S$ a morphism. Assume
\begin{enumerate}[\indent a)]
  \item the morphism $f$ is locally acyclic,
  \item for every specialisation morphism $t\to \widetilde S^s$ and every
    $q\geqslant 0$, the cospecialisation arrows
    $\h^q(X_t,\dZ/n) \to \h^q(X_s,\dZ/n)$ are bijective.
\end{enumerate}

Then the canonical homomorphism $(\eR^q f_*\dZ/n)_s\to \h^q(X_s,\dZ/n)$ is
bijective for all $q\geqslant 0$.
\end{theorem}

To prove the theorem, it is clear that one may assume
$S=\widetilde S^s$. We shall show that, for every sheaf of
$\dZ/n\dZ$-modules $\sF$ on $S$, the canonical homomorphism
$\varphi^q(\sF):(\eR^q f_* f^* \sF)_s \to \h^q(X_s,f^*\sF)$ is bijective.

Every sheaf of $\dZ/n\dZ$-modules is a filtered inductive limit of
constructible sheaves of $\dZ/n\dZ$-modules (\ref{I:4-3-3}). Moreover every
constructible sheaf of $\dZ/n\dZ$-modules embeds into a sheaf of the form
$\prod i_{\lambda *} C_\lambda$, where $i_\lambda:t_\lambda\to S$ is a
finite family of generisations of $s$ and $C_\lambda$ a free
$\dZ/n\dZ$-module of finite rank over $t_\lambda$. By the definition of the
cospecialisation arrows, condition b) means that the homomorphisms
$\varphi^q(\sF)$ are bijective if $\sF$ is of this form.

One concludes with the help of a variant of Lemma (\ref{I:4-3-6}):





\begin{lemma}\label{I:5-1-8}
Let $\cC$ be an abelian category in which filtered inductive limits
exist. Let $\varphi^\bullet:T^\bullet\to{T'}^\bullet$ be a
morphism of cohomological functors commuting with filtered inductive
limits, defined on $\cC$ with values in the category of abelian
groups. Suppose there exist two subsets $\cD$ and $\cE$ of objects
of $\cC$ such that:
\begin{enumerate}[\indent a)]
  \item every object of $\cC$ is a filtered inductive limit of objects belonging
    to $\cD$,
  \item every object belonging to $\cD$ is contained in an object belonging
    to $\cE$.
\end{enumerate}

Then the following conditions are equivalent:
\begin{enumerate}[\indent (i)]
  \item $\varphi^q(A)$ is bijective for all $q\geqslant 0$ and all
    $A\in\ob\cC$.
  \item $\varphi^q(M)$ is bijective for all $q\geqslant 0$ and all $M\in\cE$.
\end{enumerate}
\end{lemma}

The proof of the lemma is by passage to the inductive limit,
induction on $q$ and repeated application of the five lemma to the
diagram of cohomology exact sequences deduced from an exact sequence
$0\to A\to M\to A'\to 0$, with $A\in\cD$, $M\in\cE$, $A'\in \ob\cC$. 





\begin{corollary}\label{I:5-1-9}
Let $S$ be the spectrum of a strictly henselian noetherian local ring and
$f:X\to S$ a locally acyclic morphism. Suppose that, for every geometric
point $t$ of $S$, we have $\h^0(X_t,\dZ/n) = \dZ/n$ and
$\h^q(X_t,\dZ/n) = 0$ for $q>0$ (in other words the geometric fibres of
$f$ are acyclic). Then $f_* \dZ/n = \dZ/n$ and $\eR^qf_*\dZ/n = 0$ for
$q>0$.
\end{corollary}





\begin{corollary}\label{I:5-1-10}
Let $f:X\to Y$ and $g:Y\to Z$ be morphisms of locally
noetherian schemes. Then, if $f$ and $g$ are locally acyclic, so is
$g\circ f$.
\end{corollary}

One may assume that $X$, $Y$ and $Z$ are strictly local and that $f$ and $g$
are local morphisms. It is then a matter of showing that, if $z$ is an
algebraic geometric point of $Z$, we have $\h^0(X_z,\dZ/n) = \dZ/n$ and
$\h^q(X_z,\dZ/n) = 0$ for $q>0$.

Since $g$ is locally acyclic, we have $\h^0(Y_z,\dZ/n) = \dZ/n$, and
$\h^q(Y_z,\dZ/n) = 0$ for $q>0$. Moreover the morphism $f_z:X_z\to Y_z$ is
locally acyclic (\ref{I:5-1-4}) and its geometric fibres are
acyclic, since they are vanishing cycle varieties of $f$.
By (\ref{I:5-1-9}), we thus have $\eR^q {f_z}_*\dZ/n = 0$ for $q>0$. Moreover
${f_z}_*\dZ/n$ is constant with fibre $\dZ/n$ over $Y_z$. One concludes with
the Leray spectral sequence of $f_z$. 










\subsection{Local acyclicity of a smooth morphism}\label{I:5-2}





\begin{theorem}\label{I:5-2-1}
A smooth morphism is locally acyclic.
\end{theorem}

Let $f:X\to S$ be a smooth morphism. The assertion is local for the \'etale
topology on $X$ and $S$, so we may assume that $X$ is the affine space of
type of dimension $d$ over $S$. By passage to the limit, we may assume that $S$
is noetherian, and transitivity of local acyclicity (\ref{I:5-1-10}) shows
that it suffices to treat the case $d=1$.

Let $s$ be a geometric point of $S$ and $x$ a geometric point of
$X$ centred at a closed point of $X_s$. We must show that the geometric
fibres of the morphism $\widetilde X^x\to \widetilde S^s$ are
acyclic. We henceforth put $S=\widetilde S^s=\spec(A)$ and
$X=\widetilde X^x$. We have $X\simeq \spec A\{T\}$, where $A\{T\}$ is the
henselisation of $A[T]$ at the point $T=0$ above $s$.

If $t$ is a geometric point of $S$, the fibre $X_t$ is a projective limit
of smooth affine curves over $t$. Hence $\h^q(X_t,\dZ/n) = 0$
for $q\geqslant 2$ and it suffices to show that $\h^0(X_t,\dZ/n) = \dZ/n$ and
$\h^1(X_t,\dZ/n) = 0$ for $n$ prime to the residue characteristic of
$S$. This follows from the next two propositions. 





\begin{proposition}\label{I:5-2-2}
Let $A$ be a strictly henselian local ring, $S=\spec(A)$ and
$X=\spec A\{T\}$. Then the geometric fibres of $X\to S$ are connected.
\end{proposition}

One can reduce by passage to the limit to the case where $A$ is a strict
henselisation of a $\dZ$-algebra of finite type.

Let $\bar t$ be a geometric point of $S$, localised at $t$, and $k'$ a
finite separable extension of $k(t)$ in $k(\bar t)$. Put $t'=\spec(k')$
and $X_{t'} = X\times_S\spec(k')$. We must check that, whatever
$\bar t$ and $t'$ are, $X_{t'}$ is connected (by which we mean connected and
non-empty). Let $A'$ be the normalisation of $A$ in $k'$, i.e.\ the ring of
elements of $k'$ integral over the image of $A$ in $k(t)$. We have
$A\{t\}\otimes_A A'\iso A'\{T\}$: the left-hand side is indeed henselian
local (since $A'$ is finite over $A$, and local) and a limit of local
algebras \'etale over $A'\{T\}=A\{T\}\otimes_A A'$. The scheme $X_{t'}$ is thus
still the fibre at $t'$ of $X=\spec\left(A'\{T\}\right)$ over $S'=\spec(A')$.
The local scheme $X'$ is normal, hence integral; its localisation $X_{t'}$ is
still integral, a fortiori connected. 





\begin{proposition}\label{I:5-2-3}
Let $A$ be a strictly henselian local ring, $S=\spec(A)$ and
$X=\spec\left(A\{T\}\right)$. Let $\bar t$ be a geometric point of $S$
and $X_{\bar t}$ the corresponding geometric fibre. Then every Galois
\'etale covering of $X_{\bar t}$ of order prime to the characteristic of the
residue field of $A$ is trivial.
\end{proposition}





\begin{lemma}[Zariski-Nagata purity theorem in dimension $2$] \label{I:5-2-4} % originally 5.2.3.1
Let $C$ be a regular local ring of dimension $2$ and $C'$ a finite normal
$C$-algebra, \'etale over the open complement of the closed point
of $\spec(C)$. Then $C'$ is \'etale over $C$.
\end{lemma}

Indeed $C'$ is normal of dimension $2$, hence $\operatorname{prof}(C')=2$.
Since $\operatorname{prof}(C')+\dim\operatorname{proj}(C') = \dim(C)=2$, we
conclude that $C'$ is free over $C$. Then the set of points of $C$ where
$C'$ is ramified is defined by one equation, the discriminant; since it
contains no point of height $1$, it is empty. 





\begin{lemma}[special case of Abhyankar's lemma]\label{I:5-2-5} % originally 5.2.3.2
Let $S=\spec(V)$ be a trait, $\pi$ a uniformiser, $\eta$ the generic
point of $S$, $X$ smooth over $S$, irreducible, of relative
dimension $1$, $\widetilde X_\eta$ a Galois \'etale covering of $X_\eta$,
of degree $n$ invertible on $S$, and $S_1=\spec\left(V[\pi^{1/n}]\right)$.
Denote by a subscript $(-)_1$ base change from $S$ to $S_1$. Then
$\widetilde X_{1\eta}$ extends to an \'etale covering of $X_1$.
\end{lemma}

Let $\widetilde X_1$ be the normalisation of $X_1$ in $\widetilde X_{1\eta}$.
In view of the structure of the tame inertia groups of the discrete
valuation rings localised from $X$ at the generic points of the special fibre
$X_s$, $\widetilde X_1$ is \'etale over $X_1$ over the generic fibre, and at
the generic points of the special fibre. By (\ref{I:5-2-4}), it is \'etale
everywhere. 





\subsubsection{}\label{I:5-2-6} % originally 5.2.3.3

Denote by $t$ the point at which $\bar t$ is localised. We may replace
$A$ by its normalisation in a finite separable extension
$k(t')$ of $k(t)$ in $k(\bar t)$ (cf. \ref{I:5-2-2}). This, and a preliminary
passage to the limit, allow us to assume that
\begin{enumerate}[\indent a)]
  \item $A$ is normal noetherian, and $t$ is the generic point of $S$.
  \item The \'etale covering of $X_{\bar t}$ under consideration comes from an
    \'etale covering of $X_t$.
  \item It comes from an \'etale covering $\widetilde X_U$ of the inverse
    image $X_U$ of a non-empty open subset $U$ of $S$ (follows from b): $t$
    is the limit of the $U$'s).
  \item The complement of $U$ is of codimension $\geqslant 2$ (at the cost
    of enlarging $k(t)$, by applying (\ref{I:5-2-5}) to the discrete
    valuation rings localised from $S$ at the points of $S\setminus U$ of
    codimension $1$ in $S$; these points are finite in number).
  \item The covering $\widetilde X_U$ is trivial above the subscheme
    $T=0$. (At the cost of further enlarging $k(t)$).
\end{enumerate}

When these conditions are satisfied, we shall see that the covering
$\widetilde X_U$ is trivial. 





\begin{lemma}\label{I:5-2-7} % originally 5.2.3.4
Let $A$ be a normal strictly henselian noetherian local ring, $U$
an open subset of $\spec(A)$ whose complement is of codimension $\geqslant 2$,
$V$ its inverse image in $X=\spec\left(A\{T\}\right)$ and $V'$ an
\'etale covering of $V$. If $V'$ is trivial above $T=0$, then $V'$
is trivial.
\end{lemma}

Let $B=\Gamma(V',\sO)$. Since $V'$ is the inverse image of $V$ in
$\spec(B)$, it suffices to see that $B$ is finite \'etale over $A\{T\}$
(hence split, since $A\{T\}$ is strictly henselian). Let
$\widehat X=A\llbracket T\rrbracket$, and denote by $(-)^\wedge$ base change
from $X$ to $\widehat X$. The scheme $\widehat X$ is faithfully flat
over $X$. Hence
$\Gamma(\widehat V',\sO) = B\otimes_{A\{T\}} A\llbracket T\rrbracket$, and it
suffices to see that this ring $\widehat B$ is finite \'etale over
$A\llbracket T\rrbracket$.

Let $V_m$ (resp. $V_m'$) be the subscheme of $\widehat V$ (resp. $\widehat V'$)
with equation $T^{m+1} = 0$. By hypothesis, $V_0'$ is a trivial covering
of $V_0$: a sum of $n$ copies of $V_0$. Likewise for $V_m'/V_m$, since
\'etale coverings are insensitive to nilpotents. One deduces
\[
  \varphi : \Gamma(\widehat V',\sO) \to \varprojlim_m \Gamma(V_m',\sO) = \left(\varprojlim_m \Gamma(V,\sO)\right)^n \text{.}
\]
By hypothesis, the complement of $U$ is of depth $\geqslant 2$: we have
$\Gamma(V_m,\sO)=A[T]/(T^{m+1})$, and $\varphi$ is a homomorphism from
$\widehat B$ to $A\llbracket T\rrbracket^n$. Over $U$, it gives $n$
distinct sections of $\widehat V'/\widehat V$: $\widehat V'$ is therefore
trivial, a sum of $n$ copies of $\widehat V$. Since the complement of
$\widehat V$ in $\widehat X$ is still of codimension $\geqslant 2$ (hence of
depth $\geqslant 2$), we deduce that $\widehat B=A\llbracket T\rrbracket^n$,
whence the lemma. 










\subsection{Applications}\label{I:5-3}





\begin{theorem}[specialisation of cohomology groups]\label{I:5-3-1}
Let $f:X\to S$ be a proper locally acyclic morphism, for example a
proper smooth morphism. Then the sheaves $\eR^q f_*\dZ/n$ are locally
constant constructible and for every specialisation arrow
$t\to\widetilde S^s$, the cospecialisation arrows
$\h^q(X_t,\dZ/n)\to\h^q(X_s,\dZ/n)$ are bijective.
\end{theorem}

This follows immediately from the definition of the
cospecialisation arrows and the finiteness and base change theorems for
proper morphisms. 





\begin{theorem}[base change by a smooth morphism]\label{I:5-3-2}
Let there be a cartesian diagram
\[\xymatrix{
  X' \ar[r]^-{g'} \ar[d]^-{f'}
    & X \ar[d]^-f \\
  S' \ar[r]^-{g}
    & S
}\]
with $g$ smooth. For every sheaf $\sF$ on $X$, torsion and prime to the
residue characteristics of $S$, we have
\[
  g^* \eR^q f_* \sF \iso \eR^q f_*'({g'}^* \sF) \text{.}
\]
\end{theorem}

Taking an open covering of $X$, we reduce to the case where $X$ is
affine, then by passage to the limit to the case where $X$ is of finite type
over $S$. Then $f$ factors as an open immersion $j:X\to \bar X$ and a
proper morphism $\bar f:\bar X\to S$. From the Leray spectral sequence for
$\bar f\circ j$ and the base change theorem for proper morphisms,
we deduce that it suffices to prove the theorem in the case
where $X\to S$ is an open immersion.

In this case, if $\sF$ is of the form $\varepsilon_* C$, where
$\varepsilon:t\to X$ is a geometric point of $X$, the theorem is a
corollary of (\ref{I:5-1-5}). The general case follows from Lemma
(\ref{I:5-1-8}). 





\begin{corollary}\label{I:5-3-3}
Let $K/k$ be an extension of separably closed fields, $X$ a $k$-scheme
and $n$ an integer prime to the characteristic of $k$. Then the canonical
map $\h^q(X,\dZ/n) \to \h^q(X_K,\dZ/n)$ is bijective for all
$q\geqslant 0$.
\end{corollary}

It suffices to note that $\bar K$ is an inductive limit of smooth
$\bar k$-algebras. 





\begin{theorem}[relative purity]\label{I:5-3-4}
Let there be a commutative diagram
\begin{equation*}\tag{5.3.4.1}\label{I:eq:5-3-4-1}\xymatrix{
  U \ar@{^{(}->}[r]^-j \ar[dr]
    & X \ar[d]^-f
    & \ar@{_{(}->}[l]_-i Y \ar[dl]^-h \\
  & S
}\end{equation*}
with $f$ smooth purely of relative dimension $N$, $h$ smooth purely of
relative dimension $N-1$, $i$ a closed embedding and $U=X\setminus Y$. For
$n$ prime to the residue characteristics of $S$, we have
\begin{equation*}
\begin{aligned}
  j_* \dZ/n     &= \dZ/n \\
  \eR^1 j_*\dZ/n &= \dZ/n(-1)_Y \\
  \eR^q j_*\dZ/n &= 0 \qquad \text{for $q\geqslant 2$}
\end{aligned}
\end{equation*}
\end{theorem}

In these formulas, $\dZ/n(-1)$ denotes the $\dZ/n$-dual of $\dmu_n$. If $t$
is a local equation for $Y$, the isomorphism
$\eR^1j_*\dZ/n\simeq\dZ/n(-1)_Y$ is defined by the map
$a:\dZ/n\to \eR^1j_*\dmu_n$ sending $1$ to the class of the $\dmu_n$-torsor of
$n$-th roots of $t$.

The question is local in nature. This allows one to replace $(X,Y)$ by a
locally isomorphic pair, for example
\[\xymatrix{
  \dA_T^1 \ar@{^{(}->}[r]^-j \ar[dr]_-g
    & \dP_T^1 \ar[d]^-f
    & \ar@{_{(}->}[l]_-i T \ar[dl] \\
  & T
}\]
with $T=\dA_S^{n-1}$ and $i=$ the section at infinity. Corollary (\ref{I:5-1-9})
applies to $g$, and gives $\eR^q g_ \dZ/n=\dZ/n$ for $q=0$, $0$ for $q>0$.
For $f$, we have moreover \eqref{I:eq:4-6-2-1}
\[
  \eR^q f_* \dZ/n = \dZ/n,0,\dZ/n(-1),0 \text{ for $q>2$.}
\]
One checks easily that $j_* \dZ/n = \dZ/n$, and that the $\eR^q j_*\dZ/n = 0$
are concentrated on $i(T)$ for $q>0$. The Leray spectral sequence
$E_2^{pq} = \eR^p f_* \eR^q j_*\dZ/n \Rightarrow \eR^{p+q}g_* \dZ/n$ thus reduces
to
\[\xymatrix{
  i^* \eR^q j_* \dZ/n & 0 & \cdots \\
  i^* \eR^1 j_* \dZ/n \ar[drr]^-{d_2} & 0 & \cdots \\
  \dZ/n              & 0 & \dZ/n(-1) & 0 & \cdots
}\]
$\eR^q j_*\dZ/n=0$ for $q\geqslant 2$, and $\eR^1 j_*\dZ/n$ is the extension
by zero of a locally free sheaf of rank one over $T$ (isomorphic, via
$d_2$, to $\dZ/n(-1)$). The map $a$, defined above, being
injective (as one checks fibre by fibre), it is an isomorphism, and
this proves (\ref{I:5-3-4}). 





\subsubsection{}\label{I:5-3-5}

We refer to \cite[XVI.4,5]{sga4} for the proof of the following
applications of the acyclicity theorem (\ref{I:5-2-1}). 





\begin{theorem}\label{I:5-3-6} % originally 5.3.5.1
Let $f:X\to S$ be a morphism of schemes of finite type over $\dC$ and $\sF$ a
constructible sheaf on $X$. Then
\[
  \an{(\eR^q f_*\sF)} \sim \eR^q \an f_* (\an\sF)
\]
(cf. \ref{I:4-6-3}: in ordinary cohomology, it is necessary to assume that
$\sF$ is constructible and not only torsion).
\end{theorem}





\begin{theorem}\label{I:5-3-7} % originally 5.3.5.2
Let $f:X\to S$ be a morphism of schemes of finite type over a field $k$ of
characteristic $0$ and $\sF$ a constructible sheaf on $X$. Then the
$\eR^q f_* \sF$ are also constructible.
\end{theorem}

The proof uses resolution of singularities and (\ref{I:5-3-4}). It is
generalised to the case of a morphism of finite type of excellent schemes of
characteristic $0$ in \cite[XIX.5]{sga4}. Another proof,
independent of resolution, is given in this volume
(\nameref{VII}, \ref{VII:1-1})
% NOTE: originally (Th. finitude, 1.1),
it applies to a morphism of schemes of finite type over a field or over
a Dedekind ring. 




















\section{Poincar\'e duality}










\subsection{Introduction}\label{I:6-1}

Let $X$ be an oriented topological variety, purely of dimension $N$, and
suppose that $X$ admits a finite open covering
$\cU=(U_i)_{1\leqslant i\leqslant K}$, such that the non-empty intersections
of the $U_i$ are homeomorphic to balls. For such a variety,
the Poincar\'e duality theorem can be presented as follows: 

\begin{enumerate}[\indent A]
  \item The cohomology of $X$ is the \v Cech cohomology for the
    covering $\cU$. It is the cohomology of the complex
    \begin{equation}\label{I:eq:6-1}
      0 \to \dZ^{A_0} \to \dZ^{A_1} \to \cdots
    \end{equation}
    where $A_k=\{(i_0,\dotsc,i_k) : i_0<\cdots<i_k\text{ and } U_{i_0}\cap \cdots \cap U_{i_k} \ne \varnothing\}$.
    \item For $a\in A_k$, $a=(i_0,\dotsc,i_k)$, let
      $U_a=U_{i_1}\cap \cdots \cap U_{i_k}$ and let $j_a$ be the inclusion of $U_a$
      into $X$. The constant sheaf $\dZ$ on $X$ admits the (left)
      resolution
      \begin{equation}\label{I:eq:6-2}
      \xymatrix{
        \cdots \ar[r]
          & \displaystyle\bigoplus_{a\in A_1} j_{a!}\dZ \ar[r]
          & \displaystyle\bigoplus_{a\in A_0} j_{a!} \dZ \ar[r] \ar[d]
          & 0 \\
        & & \dZ \text{.}
      }
      \end{equation}
      The compactly supported cohomology $\h_c^\bullet(X,j_{a!}\dZ)$ is none
      other than the proper-support cohomology of the (oriented) ball $U_a$:
        \[
          \h_c^i(X,j_{a!}\dZ) = \begin{cases}
                                  0   & \text{if $i\ne N$} \\
                                  \dZ & \text{if $i = N$}
                                \end{cases}
        \]
\end{enumerate}

The hypercohomology spectral sequence, for the complex \eqref{I:eq:6-2}, and
cohomology with proper support, thus asserts that $\h_c^i(X)$ is the
$(i-N)$-th cohomology group of a complex
\begin{equation}\label{I:eq:6-3}
  \cdots \to \dZ^{A_1} \to \dZ^{A_0} \to 0 \text{.}
\end{equation}
This complex is the dual of the complex \eqref{I:eq:6-1}, whence Poincar\'e
duality.

The essential points of this construction are
\begin{enumerate}[\indent a)]
  \item the existence of a proper-support cohomology theory:
  \item the fact that every point $x$ of a variety $X$ purely of dimension
    $N$ has a fundamental system of open neighbourhoods $U$ for which
    \begin{equation}\label{I:eq:6-4}
      \h_c^i(U) = \begin{cases}
                    0   & \text{for $i\ne N$,} \\
                    \dZ & \text{for $i=N$.}
                  \end{cases}
    \end{equation}
\end{enumerate}

Poincar\'e duality in \'etale cohomology can be built on this
model. For $X$ smooth purely of dimension $N$ over $k$ algebraically closed,
$n$ invertible on $X$ and $x$ a closed point of $X$, the key point is to
compute the projective limit, over the \'etale neighbourhoods $U$ of $x$,
\begin{equation}\label{I:eq:6-5}
  \varprojlim \h_c^i(U,\dZ/n)
    = \begin{cases}
        0     & \text{if $i\ne 2N$} \\
        \dZ/n & \text{if $i=2N$.}
      \end{cases}
\end{equation}

Just as for topological varieties one must first treat
directly the case of an open ball (or even just that of the interval
$(0,1)$), here one must first treat directly the case of curves
(\ref{I:6-2}). The local acyclicity theorem for smooth morphisms
then reduces \eqref{I:eq:6-5} to this special case (\ref{I:6-3}).

The isomorphisms \eqref{I:eq:6-4} and \eqref{I:eq:6-5} are not canonical:
they depend on the choice of an \emph{orientation} of $X$. For $n$ invertible
on a scheme $X$, $\dmu_n$ is a sheaf of free $\dZ/n$-modules of rank
one. We denote by $\dZ/n(N)$ its $N$-th tensor power ($N\in\dZ$). The
intrinsic form of the second line of \eqref{I:eq:6-5} is
\begin{equation}\label{I:eq:6-5'}
  \varprojlim \h_c^{2N} \left(U,\dZ/n(N)\right) = \dZ/n \text{,}
\end{equation}
and $\dZ/n(N)$ is called the \emph{orientation sheaf} of $X$; since the sheaf
$\dZ/n(N)$ is constant, isomorphic to $\dZ/n$, one can take it out
of the $N$ sign and write instead
\begin{equation}\label{I:eq:6-5''}
  \varprojlim \h_c^{2N} (U,\dZ/n) = \dZ/n(-N) \text{,}
\end{equation}
and Poincar\'e duality will take the form of a perfect duality, with
values in $\dZ/n(-N)$, between $\h^i(X,\dZ/n)$ and $\h_c^{2N-i}(X,\dZ/n)$. 










\subsection{The case of curves}\label{I:6-2}





\subsubsection{}\label{I:6-2-1}

Let $\bar X$ be a smooth projective curve over $k$ algebraically closed, and
$n$ invertible on $\bar X$. The proof of (\ref{I:3-3-5}) gives, for
$\bar X$ connected, a canonical isomorphism
\[\xymatrix{
  \h^2(\bar X,\dmu_n) = \pic(\bar X)/n\pic(\bar X) \ar[r]^-{\deg}_-{\sim}
    & \dZ/n
}\]
Let $D$ be a reduced divisor on $\bar X$ and $X=\bar X\setminus D$ 
\[\xymatrix{
  X \ar@{^{(}->}[r]^-j 
    & \bar X 
    & \ar@{_{(}->}[l]_-i D \text{.}
}\]
The exact sequence $0\to j_! \dmu_n \to \dmu_n\to i_* \dmu_n\to 0$ and the fact
that $\h^i(\bar X,i_*\dmu_n) = \h^i(D,\dmu_n)=0$ for $i>0$ furnish an
isomorphism
\[\xymatrix{
  \h_c^2(X,\dmu_n) = \h^2(\bar X,j_! \dmu_n) \ar[r]^-\sim
    & \h^2(\bar X,\dmu_n) \ar[r]^-\sim
    & \dZ/n \text{.}
}\]
For $\bar X$ disconnected, we likewise have
\[
  \h_c^2(X,\dmu_n) = (\dZ/n)^{\pi_0(X)}
\]
and we define the \emph{trace morphism} as the sum
\[\xymatrix{
  \tr : \h_c^2(X,\dmu_n) \simeq (\dZ/n)^{\pi_0(X)} \ar[r]^-\Sigma
    & \dZ/n\text{.}
}\]





\begin{theorem}\label{I:6-2-2}
The pairing $\tr(a\smallsmile b)$ identifies each of the two groups
$\h^i(X,\dZ/n)$ and $\h_c^{2-i}(X,\dmu_n)$ with the dual (with values in
$\dZ/n$) of the other.
\end{theorem}

\emph{Transcendental proof.} If $\bar X$ is a smooth projective curve over
a trait $S$, and $j:X\hookrightarrow \bar X$ is the inclusion of the
complement of a divisor $D$ \'etale over $S$, the (resp.\ proper-support)
cohomologies of the special and generic geometric fibres
of $X/S$ are ``the same,'' i.e.\ the fibres of locally
constant sheaves on $S$. This follows from the analogous facts for
$\bar X$ and $D$, via the exact sequence $0\to j_! \dZ/n\to\dZ/n\to {\dZ/n}_D\to 0$
(for proper-support cohomology) and the formulas $j_*\dZ/n = \dZ/n$,
$\eR^1 j_* \dZ/n \simeq {\dZ/n}_D(-1)$, $\eR^i j_*\dZ/n = 0$ ($i\geqslant 2$)
(for ordinary cohomology) (\ref{I:5-3-4}).

This specialisation principle reduces (\ref{I:6-2-2}) to the case where $k$ is
of characteristic $0$. By (\ref{I:5-3-3}), this case reduces to the one where
$k=\dC$. Finally, for $k=\dC$, the groups $\h^\bullet(X,\dZ/n)$ and
$\h_c^\bullet(X,\dmu_n)$ coincide with the groups of the same name, computed
for the classical topological space $X_{\textnormal{cl}}$ and, via
the isomorphism $\dZ/n\to \dmu_n: x\mapsto \exp\left(\frac{2\pi i x}{n}\right)$,
the trace morphism is identified with ``integration over the fundamental
class,'' so that (\ref{I:6-2-2}) follows from Poincar\'e duality for
$X_{\textnormal{cl}}$. 





\subsubsection{Algebraic proof}\label{I:6-2-3}

For a very economical proof, see (\nameref{V}, \S\ref{V:2}).
% originally (Dualité \S 2)
% NOTE: should be hyperlinked V:2
Here is another one, related to the self-duality of the Jacobian.

Let us resume the notation of (\ref{I:6-2-1}). We may suppose -- and we
suppose -- that $X$ is connected. The cases $i=0$ and $i=2$ being trivial, we
also suppose that $i=1$. Define $_D\dG_m$ by the exact sequence
$0\to _D\dG_m \to \dG_m \to i_*\dG_m\to 0$ (sections of $\dG_m$ congruent to
$1\mod D$). The group $\h^1(\bar X, _D\dG_m)$ classifies invertible
sheaves on $\bar X$ trivialised over $D$. It is the group of points of
$\pic_D(\bar X)$, an extension of $\dZ$ (the degree) by the group of points
of a generalised Rosenlicht Jacobian (with conductor
$1$ at each point of $D$) $\pic_D^0(\bar X)$, itself an extension of the
abelian variety $\pic^0(\bar X)$ by the torus
$\dG_m^D/(\text{$\dG_m$ diagonal})$. 

\begin{enumerate}[\indent a)]
  \item The exact sequence
    $0\to j_!\dmu_n \to _D\dG_m \xrightarrow{n} _D\dG_m\to 0$ gives an
    isomorphism
    \begin{equation*}\tag{6.2.3.1}\label{I:eq:6-2-3-1}
      \h_c^1(X,\dmu_n) = \pic_D^0(\bar X)_n \text{.}
    \end{equation*}
  \item The map which to $x\in X(k)$ associates the class of the invertible
    sheaf $\sO(x)$ on $\bar X$, trivialised by $1$ over $D$, comes from
    a morphism
    \[
      f:X \to \pic_D(\bar X) \text{.}
    \]
\end{enumerate}
For what follows, fix a base point $0$ and put $f_0(x)=f(x)-f(0)$. For
every homomorphism $v:\pic_D^0(\bar X)_n \to \dZ/n$, let
$\bar v\in \h^1\left(\pic_D^0(\bar X),\dZ/n\right)$ be the image under $v$ of
the class in $\h^1\left(\pic_D^0(\bar X),\pic_D^0(\bar X)_n\right)$ of the
torsor defined by the extension
\[\xymatrix{
  0 \ar[r]
    & \pic_D^0(\bar X)_n \ar[r]
    & \pic_D^0(\bar X) \ar[r]^-n
    & \pic_D^0(\bar X) \ar[r]
    & 0
}\]
Geometric class field theory (as expounded in
Serre \cite{se59}) shows that the map $v\mapsto f_0^*(\bar v)$:
\begin{equation*}\tag{6.2.3.2}\label{I:eq:6-2-3-2}
  \hom\left(\pic_D^0(\bar X)_n,\dZ/n\right) \to \h^1(X,\dZ/n)
\end{equation*}
is an isomorphism. To deduce (\ref{I:6-2-2}) from \eqref{I:eq:6-2-3-1} and
\eqref{I:eq:6-2-3-2}, it remains to know that
\begin{equation*}\tag{6.2.2.3}\label{I:eq:6-2-3-3}
  \tr\left(u \smallsmile f_0^*(\bar v)\right) = -v(u) \text{.}
\end{equation*}

% originally (Dualité, 3.2.4).










\subsection{The general case}\label{I:6-3}

Let $X$ be a smooth algebraic variety purely of dimension $N$, over
$k$ algebraically closed. To state the Poincar\'e duality theorem,
one must first define the trace morphism
\[
  \tr : \h_c^{2N}\left(X,\dZ/n(N)\right) \to \dZ/n\text{.}
\]
The definition is a painful d\'evissage starting from the case of curves
\cite[XVIII \S 2]{sga4}. One then has 





\begin{theorem}\label{I:6-3-1}
The pairing $\tr(a\smallsmile b)$ identifies each of the groups
$\h_c^i\left(X,\dZ/n(N)\right)$ and $\h^{2 N-i}(X,\dZ/n)$ with the dual of
the other.
\end{theorem}

Let $x\in X$ be a closed point and $X_x$ the strict localisation of $X$ at
$x$. We put, for $U$ running over the \'etale neighbourhoods of $x$,
\begin{equation*}\tag{6.3.1.1}\label{I:eq:6-3-1-1}
  \h_c^\bullet(X_x,\dZ/n) = \varprojlim \h_c^\bullet(U,\dZ/n) \text{.}
\end{equation*}
It would be preferable to consider instead the pro-object
$\text{``}\varprojlim\text{''} \h_c^\bullet(U,\dZ/n)$ but, the groups at stake
being finite, the difference is inessential. As we have tried to
explain in the introduction, (\ref{I:6-3-1}) follows from the fact that
\begin{equation*}\tag{6.3.1.2}\label{I:eq:6-3-1-2}
\begin{aligned}
  \h_c^i(X_x,\dZ/n) &= 0 \text{ for $i\ne 2N$ and } \\
  \tr : \h_c^{2 N}\left(X_x,  \dZ/n(N)\right) &\iso \dZ/n \text{ is an isomorphism.}
\end{aligned}
\end{equation*}

The case $N=0$ is trivial. If $N>0$, let $Y_y$ be the strict localisation at
a closed point of a scheme $Y$ smooth purely of dimension $N-1$, and let
$f:X_x\to Y_y$ be an essentially smooth morphism (of relative dimension
one). The proof uses the Leray spectral sequence in cohomology with
proper support for $f$, to reduce to the case of curves. Since the
``proper-support cohomologies'' under consideration are defined as limits
\eqref{I:eq:6-3-1-1}, the existence of such a spectral sequence raises various
passage-to-the-limit problems, treated in all too great detail in
\cite[XVIII]{sga4}. For every geometric point $z$ of $Y_y$, we have
\[
  (\eR^if_! \dZ/n)_z = \h_c^i(f^{-1}(z),\dZ/n) \text{.}
\]
The geometric fibre $f^{-1}(z)$ is a projective limit of smooth curves
over an algebraically closed field. It satisfies Poincar\'e duality.
Its ordinary cohomology is given by the local acyclicity theorem
for smooth morphisms:
\[
  \h^i(f^{-1}(z),\dZ/n)
    = \begin{cases}
        \dZ/n & \text{for $i=0$} \\
        0     & \text{for $i>0$.}
      \end{cases}
\]
By duality, we have
\[
  \h_c^i(f^{-1}(z),\dZ/n)
    = \begin{cases}
        \dZ/n(-1) & \text{for $i=2$} \\
        0         & \text{for $i\ne 2$.}
      \end{cases}
\]
and the Leray spectral sequence reads
\[
  \h_c^i\left(X_x,\dZ/n(N)\right) = \h_c^{i-2}\left(Y_y,\dZ/n(N-1)\right) \text{.}
\]
One concludes by induction on $N$. 










\subsection{Variants and applications}\label{I:6-4}

In \'etale cohomology, one can construct a ``duality formalism''
($=$ functors $\eR f_!$, $\eR f_!$, $\eR f^!$ satisfying various
compatibilities and adjunction formulas) parallel to the one existing in
coherent cohomology. In this language, the results of the preceding
paragraph are transcribed as follows: if $f:X\to S$ is smooth purely of
relative dimension $N$, and $S=\spec(k)$ with $k$ algebraically closed,
then
\[
  \eR f^! \dZ/n = \dZ/n [2N](N) \text{.}
\]
This result holds without hypothesis on $S$. It admits the 





\begin{corollary}\label{I:6-4-1}
If $f$ is smooth purely of relative dimension $N$, and the sheaves
$\eR^i f_! \dZ/n$ are locally constant, then the sheaves
$\eR^i f_* \dZ/n$ are also locally constant, and
\[
  \eR^i f_* \dZ/n = \const\hom\left(\eR^{2 N-i} f_! \dZ/n(N),\dZ/n\right)\text{.}
\]
\end{corollary}

In particular, under the hypotheses of the corollary, the sheaves
$\eR^i f_* \dZ/n$ are constructible. Starting from this, one can show that, if
$S$ is of finite type over the spectrum of a field or a Dedekind ring,
then, for every morphism of finite type $f:X\to S$ and every constructible
sheaf $\sF$ on $X$, the sheaves $\eR^i f_* \sF$ are constructible
(\nameref{VII}, \ref{VII:1-1}). 
% originally (Th. finitude, 1.1).
